%% =============================================================================
%% FULL_THESIS.tex -- Full thesis scaffold
%% =============================================================================
%% This file is a complete LaTeX document that wraps the finished Results
%% chapter (Chapter 5) in a standard thesis structure. The other chapters
%% are scaffolds: section/subsection headings are in place with TODO markers
%% for the prose. Fill those in incrementally without disturbing Chapter 5.
%%
%% Structure:
%%   Front matter: title page, abstract, acknowledgements, TOC, LoF, LoT
%%   Chapter 1: Introduction
%%   Chapter 2: Background and Related Work
%%   Chapter 3: Problem formulation
%%   Chapter 4: Methods                       (use METHODS_BUILD_GUIDE.md)
%%   Chapter 5: Results                        (complete; do not modify body)
%%   Chapter 6: Discussion
%%   Chapter 7: Conclusion
%%   Bibliography
%%   Appendices A--E                           (per THESIS_BUILD_GUIDE sec 7)
%%
%% Compile with pdflatex + bibtex (or biber). Requires references.bib in the
%% project root.
%%
%% Figure PDFs referenced by Chapter 5 (must be on the \graphicspath):
%%   F_dataset_classes.pdf
%%   F_statistical_heterogeneity_forest.pdf
%%   F_engineered_per_class.pdf
%%   F_mu_sensitivity_outcome_and_convergence.pdf
%%   F_val_curves_extreme_gaps.pdf
%%   F_l4_four_condition_per_class_grid.pdf
%%   F_asymmetric_lr_L4.pdf
%%   F_heterogeneity_escalation.pdf
%%   F_cross_experiment_per_class.pdf
%%   F_update_norms.pdf
%% =============================================================================

\documentclass[11pt,a4paper]{report}

\usepackage[utf8]{inputenc}
\usepackage[T1]{fontenc}
\usepackage{graphicx}
\usepackage{float}   % provides [H] specifier for exact float placement (forces "here, no floating")
\graphicspath{{./}}  % flat Overleaf bundle: all figures sit next to main.tex
\usepackage{booktabs}
\usepackage{amsmath, amssymb}
\usepackage{bm}
\usepackage{microtype}
\usepackage{hyperref}
\usepackage{caption}
\usepackage{xcolor}
\usepackage[margin=2.5cm]{geometry}
\usepackage{setspace}
\usepackage[round,authoryear]{natbib}
\onehalfspacing

\title{Federated Optimisation on DermaMNIST\\
\large FedAvg, FedProx, and FedNova under Statistical and System Heterogeneity}
\author{Barbara Koch}
\date{}

\begin{document}

%% =============================================================================
%% FRONT MATTER
%% =============================================================================

\maketitle

\begin{abstract}
Federated learning lets institutions train a shared model without pooling private
data, but real federations are heterogeneous in both their data and their compute,
and medical image datasets are severely class-imbalanced. This thesis studies
federated optimisation on DermaMNIST, a seven-class dermatology benchmark
dominated by a single benign class, comparing FedAvg, FedProx, and FedNova across
controlled statistical- and system-heterogeneity regimes and reporting test
macro-F1 at the best-validation checkpoint. Rather than crowning a universal
winner, it maps the regimes in which each method preserves or loses rare-class
signal. The central finding is that under Li-style straggler asymmetry the FedProx
advantage (about $+0.115$ macro-F1) is carried almost entirely by its
$\gamma$-inexact acceptance of partial updates rather than by the proximal term,
and that FedNova is robust under deterministic learning-rate asymmetry yet
collapses under random per-round stragglers. Across the failing regimes, rare
lesions are misrouted into the majority class. These results come from a single
low-resolution benchmark with simulated heterogeneity and $n=3$ mechanism probes,
and establish no universal optimiser ranking.
\end{abstract}

\chapter*{Acknowledgements}
\addcontentsline{toc}{chapter}{Acknowledgements}
%% TODO

\tableofcontents
\listoffigures
\listoftables


%% =============================================================================
%% CHAPTER 1: INTRODUCTION
%% =============================================================================
\chapter{Introduction}
\label{ch:intro}

\section{Motivation}
\label{sec:intro-motivation}

Diagnostic models improve with large, diverse image collections, but in medicine
those collections are fragmented across institutions and cannot usually be pooled
centrally because of privacy regulation and governance. Federated learning
sidesteps this by keeping each institution's data in place and exchanging only
model updates \citep{mcmahan2017fedavg, rieke2020future, sheller2020medicalfl},
which makes it a natural setting for collaborative learning across hospitals
\citep{kairouz2021advances}. It also introduces two difficulties that are acute
in medicine. Under \emph{statistical heterogeneity}, sites see different case
mixes, so client data is non-IID; under \emph{system heterogeneity}, sites differ
in compute and availability, so some clients return only partial updates in a
round. The two coincide most damagingly when the under-resourced site is also the
one holding the rare cases.

Medical data is, in addition, severely class-imbalanced. DermaMNIST
\citep{yang2023medmnistv2}, derived from the HAM10000 dermatoscopic collection
\citep{tschandl2018ham10000} and used throughout this thesis, is dominated by
benign melanocytic nevi, while the diagnostically urgent lesions are rare. A model
can therefore post high accuracy while failing on exactly the lesions that matter,
so we report macro-F1 and per-class rare-class diagnostics rather than accuracy.
Crucially, the information needed to recognise a rare lesion reaches the global
model only through the updates of the client that holds it; if heterogeneity
under-weights or discards those updates, the rare-class signal is lost. The
quantity that governs rare-class performance is thus the \emph{flow of
rare-client signal into the global model}, and the federated optimiser --- not the
network architecture --- is the lever that protects or forfeits it.

\section{Research problem}
\label{sec:intro-problem}

The optimisers proposed for these conditions --- FedAvg \citep{mcmahan2017fedavg},
FedProx with its proximal term and tolerance for inexact updates
\citep{li2020fedprox}, and FedNova with its work-normalised aggregation
\citep{wang2020fednova} --- are usually compared by which scores higher on a task
\citep{li2022niidbench}, rarely isolating \emph{why}, often conflating the two
heterogeneities, and judging by accuracy. This thesis instead asks \emph{when}
each method preserves rare-class performance, treating the answer as a
\emph{regime map} rather than a universal ranking. It separates four levers that
are usually bundled together --- the data partition, the optimiser, the
partial-update protocol, and loss-side imbalance correction --- and varies them one
at a time. All experiments are simulated federated learning on a controlled
benchmark, so the work maps mechanisms rather than validating a clinical
deployment.

\section{Research questions and contributions}
\label{sec:intro-questions}

This framing raises four questions, examined in turn in
Chapter \ref{ch:results}. Because the literature often credits FedProx and
FedNova with broad gains, the thesis first asks whether statistical heterogeneity
alone is enough to separate FedAvg and FedProx or whether additional stress is
required; then, where FedProx does help under stragglers, whether that benefit
comes from its proximal regularisation or from its $\gamma$-inexact acceptance of
partial updates; whether FedNova's work-normalisation is robust in general or
only under deterministic learning-rate asymmetry, failing under random
round-to-round straggling; and finally whether rare-class collapse into the
majority class is the shared failure mode, and whether the FedProx advantage is
merely class-imbalance correction in disguise.

Answering these yields the contributions of the thesis. It maps the regimes in
which FedAvg, FedProx, and FedNova preserve or lose rare-class signal; it
decomposes the FedProx straggler-regime advantage into its proximal and
$\gamma$-inexact components using matched class-disjoint (L4) and quantity-only
(L1) partitions that fix when the effect can appear; it characterises FedNova's
regime-dependence across deterministic and random system heterogeneity; it
identifies rare-class collapse into the dominant melanocytic-nevi class as the
common failure mode; and it shows, through a loss-side comparison, that the
FedProx advantage is a drift-and-protocol effect rather than class-imbalance
correction in disguise.


%% =============================================================================
%% CHAPTER 4: METHODS
%% =============================================================================
\chapter{Methods}
\label{ch:methods}

\section{Problem setup and design}
\label{sec:method-problem}
\label{sec:method-rationale}

We consider a federation of $K$ clients coordinated by a central server that
together train a single global model without exchanging raw data
(Figure~\ref{fig:fl-workflow}). Client $i$ holds a private training set $D_i$ of
$n_i$ examples, with $N = \sum_i n_i$. In each round $t$ the server broadcasts the
global model $w^{t}$; every participating client runs $E$ epochs of local SGD on
$D_i$ to obtain $w_i^{t+1}$ and returns its update $d_i = w_i^{t+1} - w^{t}$; the
baseline rule, FedAvg \citep{mcmahan2017fedavg}, recombines these by the
sample-weighted mean $w^{t+1} = \sum_i (n_i / N)\, w_i^{t+1}$. FedProx
\citep{li2020fedprox} and FedNova \citep{wang2020fednova} modify, respectively,
the local objective and the way the updates $d_i$ are normalised before
combination (\S\ref{sec:method-optimisers}).

\begin{figure}[t]
\centering
\includegraphics[width=0.86\textwidth]{F_federated_learning_workflow.pdf}
\caption{Federated learning loop (conceptual; \emph{simulated} on a single
machine in this thesis). Each round the server broadcasts the global model
$w^{t}$, clients train locally on private data $D_i$ and return model updates
$d_i$ (not raw data), and the server aggregates them into $w^{t+1}$; the loop
repeats for $R$ rounds.}
\label{fig:fl-workflow}
\end{figure}

The aim is not to crown a universal winner among the three optimisers but to map
which regimes preserve or lose rare-class performance. The design separates three
axes the FL literature often conflates.
\emph{Statistical heterogeneity} controls which clients own the rare-class signal
--- from IID, through Dirichlet label skew \citep{hsu2019dirichlet}, to
class-disjoint clients --- and is fixed before training by the data partition; it
determines \emph{where} rare-class information resides. \emph{System
heterogeneity} controls whether that signal reaches aggregation: a straggler
completes fewer local steps, and its partial update may be dropped or accepted, or
it may take smaller steps under a lower learning rate. \emph{Optimiser and
protocol choices} then determine the global response. The two heterogeneities
interact most sharply when the client carrying the rare classes is also the one
that is slow or under-weighted.

This motivates the concept that organises the thesis, \emph{rare-client signal
flow}: the per-round contribution that the client holding rare-class examples
makes to the global update. A rare class is learned only to the extent that the
information in that client's update $d_i$ reaches and persists in $w^{t+1}$, so the
organising hypothesis is that rare-class performance is preserved precisely when
rare-client signal continues to reach the global model. Two matched two-client
partitions isolate this (Figure~\ref{fig:l1-l4-schematic}): the class-disjoint
\textbf{L4}, in which one client uniquely holds melanoma, dermatofibroma, and
vascular lesions, and its quantity-only control \textbf{L1}, in which both clients
carry near-global proportions --- so suppressing a client removes unique
rare-class signal under L4 but not under L1. Performance is measured throughout by
test macro-F1 at the best-validation checkpoint, with the full metric stack and
reporting protocol defined in \S\ref{sec:method-metrics}.

Several claims are deliberately \emph{not} made. The simulator does not emulate a
real hospital network: no wall-clock deadlines, no privacy or communication-cost
budgets, no clinical validation. The optimiser set is not an exhaustive FL
benchmark (no SCAFFOLD, FedDyn, or FedAdam). The four-condition decomposition
lacks a FedAvg~$+~\gamma$-inexact cell, so the attribution reported in
Chapter~\ref{ch:results} is for the \emph{FedProx~$+~\gamma$-inexact protocol},
not a proof that $\gamma$-inexact acceptance is universally beneficial as a
standalone mechanism. The chapter follows a causal-contrast sequence that the
Results chapter walks in order: IID baseline $\to$ engineered
statistical-heterogeneity $\to$ $\mu$-sweep $\to$ L1/L4 four-condition
decomposition $\to$ learning-rate-asymmetry envelope $\to$ loss-side comparators.

\section{Dataset, metrics, and reporting}
\label{sec:method-dataset}
\label{sec:method-metrics}

All experiments use DermaMNIST from MedMNIST~v2 \citep{yang2023medmnistv2}, a
seven-class dermatology benchmark derived from the HAM10000 dataset of
dermatoscopic skin-lesion images \citep{tschandl2018ham10000}. The task is
single-image classification at $28 \times 28$ RGB resolution. Splits follow the
canonical MedMNIST~v2 release with $7007$ training, $1003$ validation, and $2005$
test images, permuted to channel-first float, scaled to $[0, 1]$, and normalised
with ImageNet channel statistics.

The benchmark is severely class-imbalanced (Figure~\ref{fig:dataset-classes},
Table~\ref{tab:dataset-classes}): the
dominant melanocytic-nevi class (mel-nevi) accounts for $66.9\,\%$ of the training
set and $1341 / 2005$ test images. The two smallest classes --- dermatofibroma
($1.1\,\%$, $23$ test images) and vascular lesions ($1.4\,\%$, $29$ test) --- have
so few support that any optimiser that fails to learn them suffers a large
macro-F1 penalty while keeping raw accuracy almost unchanged. Melanoma
($11.1\,\%$, $223$ test) is mid-support but is included in the canonical
rare-class group used throughout, $\mathrm{RARE\_IDX} = \{\text{dermato},
\text{melanoma}, \text{vascular}\}$, because missing a melanoma is the
highest-priority failure the dataset can express. Mel-nevi acts throughout as the
\emph{majority attractor}: under several failing optimiser/protocol combinations,
rare-class images are disproportionately routed into it, so raw accuracy stays
misleadingly high while macro-F1 collapses. Client partitions
(\S\ref{sec:method-partitions}) apply to the training set only; the validation and
test sets remain global and class-stratified, so any reported test-set rare-class
collapse is not a partition artefact.

\begin{table}[H]
\centering
\small
\caption{DermaMNIST class supports and roles in this thesis. Train
prevalence and test support are computed directly from the MedMNIST~v2
release; the train support totals $7007$ and the test support totals
$2005$. The rare-class group $\mathrm{RARE\_IDX} = \{\text{dermato},
\text{melanoma}, \text{vascular}\}$ defines the rare-class F1 used
throughout Chapter~\ref{ch:results}.}
\label{tab:dataset-classes}
\resizebox{\textwidth}{!}{%
\begin{tabular}{l c c p{0.42\linewidth}}
\hline
Class                              & Train prevalence & Test support & Role in thesis \\
\hline
mel-nevi (melanocytic nevi)        & $66.9\,\%$       & $1341$       & Majority attractor; collapse-mode prediction class \\
melanoma                           & $11.1\,\%$       & $223$        & Mid-support, high clinical priority; in rare-class group \\
dermatofibroma                     & $1.1\,\%$        & $23$         & Smallest class; in rare-class group \\
vascular lesions                   & $1.4\,\%$        & $29$         & Second smallest; in rare-class group \\
remaining (akiec, bcc, b-kerat)    & $19.4\,\%$       & $389$ total  & Per-class reporting only; no special role \\
\hline
\end{tabular}%
}
\end{table}

\begin{figure}[H]
\centering
\includegraphics[width=0.7\textwidth]{F_dataset_classes.pdf}
\caption{DermaMNIST class distribution. Mel-nevi dominates at $66.9\,\%$ of
the training set; dermato ($1.1\,\%$), vascular ($1.4\,\%$), and melanoma
(${\approx}11\,\%$) define the rare-class set used throughout the experiments.}
\label{fig:dataset-classes}
\end{figure}

The primary metric throughout is \emph{test macro-F1 at the best-validation
checkpoint}, computed on the global test set: macro-F1 weights all seven classes
equally and is sensitive to rare-class collapse, whereas raw accuracy saturates at
the $\approx 66.9\,\%$ mel-nevi floor and would not separate a model that has lost
the rare classes from one that has not. Secondary diagnostics are per-class F1;
per-class recall (the confusion-matrix diagonal $P(\hat{y} = c \mid y = c)$);
rare-class F1 over $\mathrm{RARE\_IDX}$ (the $0.265$ value from earlier drafts that
excluded melanoma is not used); rare-to-mel-nevi misrouting, the mean over rare
classes of $P(\hat{y} = \text{nv} \mid y = c)$; convergence curves; and per-client
update norms (diagnostic only). Balanced accuracy is stored only as a sanity check
and never used to rank optimisers --- raw accuracy would mis-rank the L4 5{:}1
LR-asymmetry cell (\S\ref{sec:lr-asymmetry}). Validation macro-F1 selects the best
checkpoint round-by-round; test metrics are evaluated once at that checkpoint.
Probabilities are not stored, so AUC and calibration are not used. A
\emph{protocol flip} is a sign reversal of $\Delta = \mathrm{FedProx} -
\mathrm{FedAvg}$ between the best-validation, final-round, and last-$K$ summaries;
it is a robustness check on reporting choice, not a result claim.

Three seed tiers are reported distinctly (Table~\ref{tab:seed-tiers}); sample SDs
use \texttt{ddof=1}. The \emph{node-pinned L4} configuration is a variance control,
not a partition design: it re-runs L4 cells on a single compute node with matched
seeds to isolate cross-node CUDA non-determinism from any optimiser effect.

\begin{table}[H]
\centering
\small
\caption{Seed-tier reporting protocol. Wins is the number of paired
within-seed comparisons in which FedProx exceeds FedAvg. Wilcoxon refers
to the paired signed-rank test, used only where $n = 10$ provides
sufficient power.}
\label{tab:seed-tiers}
\begin{tabular}{c p{0.32\linewidth} p{0.36\linewidth}}
\hline
Tier         & Statistical reporting                                                                            & Use case                                                                                                  \\
\hline
$n = 10$     & Mean $\pm$ sample SD; paired $\Delta$; wins; paired Wilcoxon signed-rank where meaningful        & Headline statistical-heterogeneity, $\mu$-sweep, IID baseline, system-het 7-client experiments            \\
$n = 3$      & Mean $\pm$ sample SD; directional only; no significance language                                & Mechanism/asymmetry probes: four-condition, perfect-storm, LR asymmetry, weighted-CE, asymmetric-$\mu$    \\
$n = 1$      & Labelled \emph{pilot}; reported only for orientation                                            & Heterogeneity ladder L0--L4; never used to derive headline numbers                                        \\
\hline
\end{tabular}
\end{table}

\section{Simulation and model}
\label{sec:method-simulation}

All experiments are simulated federated learning, not a real-hospital deployment.
Each run executes the loop of Figure~\ref{fig:fl-workflow} for $R = 150$ rounds
with $E = 20$ local epochs of mini-batch SGD per client, using the same global,
class-stratified validation and test sets throughout --- only the client training
partition varies. Two federation sizes are used. The \textbf{7-client} federation
carries the headline 10-seed benchmarks (IID, engineered balanced-paired,
Dirichlet $\alpha = 0.1$, specialist 1-of-7) and the random-straggler system-het
comparisons including the FedNova random-$\tau$ collapse; the \textbf{2-client}
federation carries the controlled mechanism probes that need a clean rare-client
vs dominant-client interaction (L1/L4 four-condition, perfect-storm, LR-asymmetry
envelope, weighted-CE/focal comparators, asymmetric $\mu$, node-pinned L4, ladder
pilot). No cross-scale generalisation between the two sizes is claimed. Two
training regimes are used: \textbf{Regime~A} (7-client default: SGD momentum
$0.9$, batch $32$, $\textrm{lr} = 0.01$, weight decay $0$) and \textbf{Regime~B}
(Li-style plain SGD for the 2-client probes: momentum $0$, batch $10$, same
lr/wd/$E$/$R$); the four-condition decomposition additionally fixes Client~1's
$E_{\textrm{straggler}} = 5$. Each experiment's regime is stated in the
corresponding Results section.

The model is a small CNN used as the benchmark vehicle, not a contribution: four
convolutional blocks of widths $32 \to 64 \to 128 \to 256$ (each a $3 \times 3$
convolution with GroupNorm \citep{wu2018groupnorm} and ReLU; max-pooling between
the first three, adaptive average pooling after the fourth), then a $256 \to 128$
dense layer with dropout $0.2$ and a seven-class head. GroupNorm is used rather
than BatchNorm because BN's per-batch statistics diverge across heterogeneous
clients, a documented non-IID FL pathology \citep{hsieh2020quagmire}. The local
objective is unweighted cross-entropy unless \S\ref{sec:method-losses} overrides,
with best-validation checkpointing.

\label{sec:method-runtimes}
Two runtimes are used: a \emph{pure-PyTorch reference loop} and the \emph{Flower
simulation} (framework $1.23$ \citep{beutel2020flower}), which carries every
experiment and all reported headline numbers. Where the two disagree on magnitude
(the engineered partition, \S\ref{sec:statistical-heterogeneity}), Flower is the
central runtime and the pure-PyTorch estimate is direction-only corroboration.
Full implementation provenance, saved-artefact formats, and reproducibility guards
are given in Appendix~\ref{app:provenance}.

\section{Partitions, optimisers, and protocols}
\label{sec:method-partitions}

Partition design is the primary lever for statistical heterogeneity: it determines
whether each client sees a near-global class mix, or rare-class signal is
concentrated on a few clients --- or, in the strongest probes, on a single client
whose contribution can then be made to fail or succeed by varying the
system-heterogeneity protocol. The partitions and their downstream consumers are
listed in Table~\ref{tab:partition-designs}. The \emph{IID} partition is a
negative control: every client receives a representative random sample, so a
substantial FedProx advantage here would not support the
heterogeneity-correction interpretation and would require implementation scrutiny.
The \emph{engineered balanced-paired} 7-client partition is a controlled
middle ground between IID and the singleton specialist case: every minority class
is concentrated on exactly two clients (rare-class signal unevenly shared) but no
minority class is isolated to a single client. It is the substrate for the
headline FedAvg/FedProx comparison and the $\mu$-sweep
(Figure~\ref{fig:engineered-partition}); by the author's own design note it was
\emph{``designed to maximise FedProx's advantage on the per-class drift
mechanism''}, so the small Flower-paired $\Delta$ reported in
\S\ref{sec:statistical-heterogeneity} should be read against that FedProx-favouring
framing rather than as a naturalistic deployment.

\begin{figure}[H]
\centering
\includegraphics[width=0.92\textwidth]{F_engineered_partition_heatmap.pdf}
\caption{Engineered balanced-paired 7-client partition: training-sample
counts per (class, client) cell. Rare classes (df, mel, vasc; outlined
in red) sit on exactly two clients each; mel-nevi (majority) is broadly
distributed across all seven clients in $\approx 670$-image shares
plus a $673$-image generalist on C6. This creates controlled
statistical heterogeneity --- rare-class signal is unevenly shared ---
without reducing any rare class to a singleton specialist case. The
bottom bar gives the total training-sample count per client.}
\label{fig:engineered-partition}
\end{figure}

The two-client \textbf{L1} and \textbf{L4} partitions are matched mechanism
controls (Figure~\ref{fig:l1-l4-schematic}): both use the same $86/14$ sample
split and protocol family and differ in one variable --- whether Client~1 carries
class information Client~0 lacks. In \emph{L4} the split is class-disjoint:
Client~0 holds mel-nevi, benign-keratosis, actinic keratoses, and basal cell
carcinoma, while Client~1 uniquely holds melanoma, dermatofibroma, and vascular
lesions (JS divergence $1.0$; the filename
\texttt{two\_client\_90\_10\_rare\_stress} is a legacy of the $90/10$ label). L4 is
the substrate for the four-condition decomposition, perfect-storm, asymmetric-LR,
and loss-side experiments. \emph{L1} keeps the same $86/14$ quantity skew but
stratifies every class across both clients in near-global proportions (JS
divergence $\approx 5 \times 10^{-6}$), so it tests whether $\gamma$-inexact
acceptance helps generically, while L4 tests whether it helps specifically when
the under-trained client holds unique rare-class signal.

\begin{figure}[H]
\centering
\includegraphics[width=\textwidth]{F_l1_l4_partition_schematic.pdf}
\caption{Two-client mechanism partitions. \textbf{L1} keeps the (86/14)
quantity split but gives both clients near-global class proportions.
\textbf{L4} is class-disjoint: Client~0 holds the common classes and
Client~1 uniquely holds melanoma, dermatofibroma, and vascular lesions.
The contrast isolates whether the under-trained client carries unique
rare-class signal.}
\label{fig:l1-l4-schematic}
\end{figure}

\begin{table}[H]
\centering
\small
\caption{Client-partition designs used in the chapter. The engineered
balanced-paired partition and the L1/L4 pair are load-bearing; the remaining
partitions test robustness or provide pilot context.}
\label{tab:partition-designs}
\resizebox{\textwidth}{!}{%
\begin{tabular}{l c p{0.22\linewidth} p{0.25\linewidth} l}
\hline
Partition                       & Federation & Statistical-het type            & Purpose                                                                & Results §                                       \\
\hline
IID                             & 7-client   & none                            & Negative control for FedProx under no statistical heterogeneity        & §\ref{sec:statistical-heterogeneity}            \\
Engineered balanced-paired      & 7-client   & paired minority ownership       & Main FedAvg/FedProx comparison and $\mu$-sweep                         & §\ref{sec:statistical-heterogeneity}, §\ref{sec:mu-sensitivity} \\
Dirichlet $\alpha = 0.1$        & 7-client   & stochastic label skew           & Test whether the engineered result generalises                         & §\ref{sec:statistical-heterogeneity}            \\
Specialist 1-of-7               & 7-client   & singleton minority ownership    & Test stronger class specialisation without changing client sizes       & §\ref{sec:statistical-heterogeneity}            \\
Ladder L0--L4                   & 2-client   & escalating label skew           & Single-seed pilot; orientation only                                    & §\ref{sec:statistical-heterogeneity} (pilot)    \\
L1                              & 2-client   & quantity-only                   & Negative control for the $\gamma$-inexact mechanism                    & §\ref{sec:li-straggler}                         \\
L4                              & 2-client   & class-disjoint                  & Mechanism substrate for straggler, LR-asymmetry, and loss experiments  & §\ref{sec:li-straggler}, §\ref{sec:lr-asymmetry}, §\ref{sec:rare-class-collapse} \\
\hline
\end{tabular}%
}
\end{table}

\subsection{Optimisers}
\label{sec:method-optimisers}

Optimiser choice and straggler-handling protocol are treated as \emph{independent}
axes: the optimisers below define how each client's local update is produced and
combined, while drop versus $\gamma$-inexact handling separately determines whether
a straggler's partial update enters aggregation. This separation is what lets the
four-condition decomposition attribute the rare-class-rescue effect to one
mechanism rather than the other (Table~\ref{tab:optimisers}). \textbf{FedAvg}
\citep{mcmahan2017fedavg} is the baseline: clients run $E$ epochs of local SGD and
the server takes the sample-weighted mean, with no drift control or work
normalisation, so on L4 the rare-class client ($\approx 14\,\%$ of samples) is
structurally under-weighted. \textbf{FedProx} \citep{li2020fedprox} keeps that
aggregation rule but augments each client's local objective with a proximal
penalty
\[
  \mathcal{L}_{\textrm{prox}}(w; w_{\textrm{global}})
  \;=\; \mathcal{L}_{\textrm{CE}}(w) \;+\; \tfrac{\mu}{2}\,
        \lVert w - w_{\textrm{global}} \rVert^{2},
\]
where $w_{\textrm{global}}$ is the round-start model and $\mu$ the proximal
coefficient (default $\mu = 0.01$ outside the $\mu$-sweep of
\S\ref{sec:mu-sensitivity}). By construction $\mu = 0$ makes the update
bit-identical to FedAvg, an equivalence verified numerically by a regression test,
so any later FedProx advantage cannot be a coding artefact. \textbf{FedNova}
\citep{wang2020fednova} keeps FedAvg-style solvers but normalises each client's
update by its effective local work $\tau_i$ via the momentum-aware factor
$a_i = (\tau_i(1-m) - m(1-m^{\tau_i}))/(1-m)^{2}$ (which reduces to $a_i = \tau_i$
when $m = 0$, as under Regime~B), aggregating $d_i / a_i$ and scaling the global
step by $a_{\textrm{eff}} = \sum_i p_i a_i$; both its asymmetry regimes lie outside
the stable-$\tau$ assumption of the original analysis.

\begin{table}[H]
\centering
\small
\caption{Federated optimisation methods used in this thesis. The
``mechanism'' column names the operative change relative to FedAvg;
``main knob'' is the hyperparameter the experimenter actually moves;
``expected vulnerability'' summarises the regime in which the method is
predicted to underperform, and the right-most column locates the
Results section in which that prediction is tested.}
\label{tab:optimisers}
\resizebox{\textwidth}{!}{%
\begin{tabular}{l p{0.20\linewidth} p{0.18\linewidth} p{0.26\linewidth} l}
\hline
Method   & Mechanism                                                & Main knob                   & Expected vulnerability                                                                  & Results §                                       \\
\hline
FedAvg   & Sample-weighted aggregation; no drift control            & none beyond local SGD       & Rare-class client structurally under-weighted under quantity skew                       & §\ref{sec:statistical-heterogeneity}, §\ref{sec:li-straggler} \\
FedProx  & FedAvg plus local proximal penalty $\tfrac{\mu}{2}\lVert w - w_{\textrm{global}}\rVert^{2}$ & $\mu$ (default $0.01$)      & $\mu$ inverted-$U$: too small $\approx$ FedAvg; too large over-constrains rare-class learning & §\ref{sec:mu-sensitivity}, §\ref{sec:li-straggler}            \\
FedNova  & Sample-weighted aggregation of $\tau_i$-normalised local directions $d_i / a_i$ & client momentum $m$, local-step count $\tau_i$ & Fails under random per-round $\tau$ (outside the stable-$\tau$ assumption of \citep{wang2020fednova}) & §\ref{sec:lr-asymmetry}                          \\
\hline
\end{tabular}%
}
\end{table}

\subsection{Partial-update protocols and stress regimes}
\label{sec:method-protocols}

System heterogeneity is simulated by reducing certain clients' local work in a
round (computational asymmetry, not wall-clock latency), via \texttt{uniform},
\texttt{fixed\_stragglers} (named clients run $E_{\textrm{straggler}} < E_{\max}$),
or \texttt{random\_stragglers} (a fraction per round run a random
$E \in [1, E_{\max}-1]$) schedules. When a client returns a partial update the
server either \emph{drops} it (excluded from that round's aggregation) or accepts
it under \emph{$\gamma$-inexact} handling (Figure~\ref{fig:signal-flow-gamma}).
\textbf{Drop versus $\gamma$-inexact is a protocol choice, not an optimiser
property:} either policy pairs with either FedAvg or FedProx, and this
independence is what the four-condition decomposition exploits.

\begin{figure}[t]
\centering
\includegraphics[width=0.82\textwidth]{F_signal_flow_gamma_inexact.pdf}
\caption{Rare-client signal flow under drop and $\gamma$-inexact partial-update
handling in the L4 partition. Client~1 uniquely holds melanoma, dermatofibroma,
and vascular examples. Under \textbf{(a)} drop handling its partial update is
discarded, so rare-class signal is lost; under \textbf{(b)} $\gamma$-inexact
handling the partial update enters aggregation and the signal is preserved. The
schematic is conceptual; the evidence in \S\ref{sec:li-straggler} concerns the
FedProx~$+$~$\gamma$-inexact protocol in this design.}
\label{fig:signal-flow-gamma}
\end{figure}

The \emph{four-condition decomposition} runs on L4 under Regime~B with the
fixed-straggler schedule ($E_{\max} = 20$, $E_{\textrm{straggler}} = 5$ on
Client~1, the rare-class holder), $n = 3$ matched seeds per cell. Holding the
partition and schedule fixed isolates the drop-vs-accept axis from any per-round
variability, so the only quantities that change between the four cells
(Table~\ref{tab:four-condition-design}) are the optimiser (FedAvg vs FedProx) and
the partial-update protocol (drop vs $\gamma$-inexact). The design deliberately
omits a FedAvg~$+~\gamma$-inexact cell, so attributing the rare-class-rescue effect
to the $\gamma$-inexact protocol (rather than to the proximal term alone) is a
claim about the \emph{FedProx~$+~\gamma$-inexact} protocol, not about
$\gamma$-inexact acceptance in isolation.

Three further system-asymmetry regimes build on this. The \emph{perfect-storm}
stress keeps L4/Regime~B but raises straggling to a random fraction $0.9$,
comparing FedAvg+drop with FedProx+$\gamma$-inexact at $\mu \in \{1.0, 0.01\}$
($n = 3$); because its schedule and arms differ from the four-condition design it
is reported as a stress demonstration of the same ordering, not a second
decomposition. The \emph{learning-rate-asymmetry} envelope keeps Client~0 at
$\textrm{lr} = 0.01$ and varies Client~1 across six ratios from $1{:}1$ to $50{:}1$
($\textrm{lr}_1 \in \{0.01, 0.005, 0.002, 0.001, 0.0005, 0.0002\}$), weakening the
rare client's update magnitude \emph{without dropping it} and so separating the
``silent client'' failure mode from the ``smaller-step'' one. The
\emph{random-$\tau$ FedNova stress} (\texttt{system\_het\_random\_fednova}) runs on
the 7-client engineered partition under Regime~A with random-straggler fraction
$0.5$, $n = 10$ --- not an L4 experiment --- deliberately making $\tau_i$ a
per-round random variable that violates the stable-$\tau$ assumption of
\citet{wang2020fednova}.

\begin{table}[H]
\centering
\small
\caption{Four-condition decomposition design (L4, two-client, Regime~B,
$E_{\textrm{straggler}} = 5$ on Client~1, $n = 3$ matched seeds per
cell). Holding the partition and schedule fixed across all four cells
isolates the optimiser-vs-protocol contribution. The design omits the
FedAvg + $\gamma$-inexact cell, so the $\gamma$-inexact attribution is
to the FedProx + $\gamma$-inexact \emph{protocol} rather than to
$\gamma$-inexact acceptance as a standalone universal mechanism.}
\label{tab:four-condition-design}
\resizebox{\textwidth}{!}{%
\begin{tabular}{l c c c p{0.28\linewidth}}
\hline
Condition                            & Straggler schedule           & Partial-update handling & Proximal term & Mechanism isolated                                          \\
\hline
FedAvg, no straggler                 & uniform ($E = 20$ everywhere) & not applicable          & --            & No-system-heterogeneity baseline                            \\
FedAvg + drop                        & fixed, Client 1 $E = 5$       & drop                    & --            & Cost of dropping the rare-class client's partial update     \\
FedProx + drop                       & fixed, Client 1 $E = 5$       & drop                    & $\mu = 0.01$  & Proximal-only contribution under matched drop handling      \\
FedProx + $\gamma$-inexact           & fixed, Client 1 $E = 5$       & accept ($\gamma$-inex)  & $\mu = 0.01$  & Additional effect of accepting the partial update           \\
\hline
\end{tabular}%
}
\end{table}

\subsection{Loss-side interventions}
\label{sec:method-losses}

The loss-side comparators test a falsification hypothesis: if FedProx's advantage
is really a disguised class-imbalance correction, a direct loss-level correction
should match or exceed it on the same partition. Three losses are compared on the
L4 substrate: plain cross-entropy (the unweighted baseline); weighted CE with
inverse-frequency weights $w_c = (1/n_c)\cdot C / \sum_j (1/n_j)$ normalised so
$\sum_c w_c = C$; and focal loss \citep{lin2017focal} with $\gamma = 2$ and a
default inverse-frequency $\alpha_c$ (so the focal comparator is focal with
class-balanced $\alpha$). Weighted CE intervenes in per-class loss contribution
during local training while FedProx intervenes in aggregation, so if FedProx were
class re-weighting in disguise, weighted CE would close the gap on its own. This is
a focused comparator set, run only on L4.

\section{Experiment families}
\label{sec:method-roadmap}

Methods has now fixed the partitions, optimisers, protocols, losses, runtime, and
seed tiers. Table~\ref{tab:experiment-families} maps the eleven experiment families
to the Results sections they support; the Results chapter follows this table as a
single argument rather than a list of detached experiments. L4 appears in five
rows --- the controlled-mechanism substrate --- while L1 appears once but is the
load-bearing negative control for the $\gamma$-inexact rare-class-rescue mechanism
(\S\ref{sec:li-straggler}).

\begin{table}[H]
\centering
\small
\caption{Experiment families and the Results sections they support. The
table bridges the design choices defined in §\ref{sec:method-problem}--§\ref{sec:method-losses}
and Results (Chapter~\ref{ch:results}). Each Results section examines
one row or a group of rows; none is presented as a detached standalone
experiment.}
\label{tab:experiment-families}
\resizebox{\textwidth}{!}{%
\begin{tabular}{l p{0.18\linewidth} p{0.22\linewidth} p{0.28\linewidth} l}
\hline
Experiment family                & Federation / partition                       & Algorithms / protocols                                  & What it tests                                                                 & Results §                                       \\
\hline
IID control                      & 7-client IID                                 & FedAvg vs FedProx, $n = 10$                              & Falsify FedProx-helps-without-heterogeneity                                  & §\ref{sec:statistical-heterogeneity}            \\
Engineered statistical-het       & 7-client balanced-paired                     & FedAvg vs FedProx, both runtimes, $n = 10$               & Headline $\Delta$ under engineered label skew                                & §\ref{sec:statistical-heterogeneity}            \\
Robustness partitions            & 7-client Dirichlet / specialist              & FedAvg vs FedProx, $n = 10$                              & Test engineered direction outside its framing                                & §\ref{sec:statistical-heterogeneity}            \\
Node-pinned L4 control           & 2-client L4, single compute node             & FedAvg vs FedProx, $n = 3$ matched                       & Isolate cross-node CUDA non-determinism                                       & §\ref{sec:statistical-heterogeneity}            \\
$\mu$ sensitivity                & 7-client balanced-paired                     & FedProx at $\mu \in \{0, 0.001, 0.01, 0.1, 1.0\}$, $n = 10$ & Locate proximal-coefficient peak; show large-$\mu$ harm                      & §\ref{sec:mu-sensitivity}                       \\
Li-style four-condition          & 2-client L4 and L1 (control)                 & FedAvg/FedProx $\times$ drop/$\gamma$-inex, $n = 3$       & Isolate proximal-only from $\gamma$-inex contribution                         & §\ref{sec:li-straggler}                         \\
Perfect-storm                    & 2-client L4, batch $10$, momentum $0$, $90\%$ random stragglers & FedAvg+drop vs FedProx+$\gamma$-inex ($\mu \in \{1.0, 0.01\}$), $n = 3$ & Stress demonstration of the four-condition ordering                        & §\ref{sec:li-straggler}                         \\
LR asymmetry                     & 2-client L4, 6 ratios $1{:}1$ to $50{:}1$    & FedAvg, FedProx, FedNova, $n = 3$ per ratio              & FedNova robustness under deterministic LR skew                                & §\ref{sec:lr-asymmetry}                         \\
FedNova random-$\tau$            & 7-client balanced-paired, $f_{\textrm{straggle}} = 0.5$ & FedNova vs system\_het\_random FA/FP, $n = 10$           & Random-$\tau$ stress outside Wang 2020's stable-$\tau$ assumption             & §\ref{sec:lr-asymmetry}                         \\
Loss-side correction             & 2-client L4                                  & FedAvg/FedProx $\times$ CE / weighted CE / focal, $n = 3$ & Test whether FedProx is a disguised class-imbalance corrector                & §\ref{sec:rare-class-collapse}                  \\
Mechanism diagnostics            & cross-protocol pool of L4 runs               & update-norm trajectories, protocol-flip audit             & Drift-control mechanism evidence; reporting-protocol robustness               & §\ref{sec:mechanism-robustness}                  \\
\hline
\end{tabular}%
}
\end{table}

\noindent{\footnotesize Due to restricted HPC access, experiments were not always
repeated with the same number of random seeds; the seed count for each experiment
is reported in the table above.}


%% =============================================================================
%% CHAPTER 5: RESULTS  --  COMPLETE; DO NOT MODIFY BODY
%% =============================================================================
%% The Results chapter is finalised. Its content is preserved verbatim below.
%% (Provenance notes pointing at the source repository were removed for this
%% flat Overleaf bundle; see the repository for source files and verification.)

\chapter{Results}
\label{ch:results}




%%%%%%%%%%%%%%%%%%%%%%%%%%%%%%%%%%%%%%%%%%%%%%%%%%%%%%%%%%%%%%%%%%%%%%%%%%%%%%%
%% sec 5.1 -- Evaluation setting and reporting protocol
%%%%%%%%%%%%%%%%%%%%%%%%%%%%%%%%%%%%%%%%%%%%%%%%%%%%%%%%%%%%%%%%%%%%%%%%%%%%%%%

% =============================================================================
% Results sec 5.1 -- Evaluation setting and reporting protocol
% Compressed. Definitional, not interpretive. Target ~250 words.
% Verified facts: class supports (66.9% / 1.1% / 1.4%); test counts;
% RARE_IDX = {dermato, melanoma, vascular} (canonical).
% =============================================================================

Every configuration is evaluated by test macro-F1 at the best-validation checkpoint
(\S\ref{sec:method-metrics}): on DermaMNIST, where melanocytic nevi account for
$66.9\,\%$ of the data, accuracy saturates at the
majority floor and hides rare-class failure, whereas macro-F1 weights all seven
classes equally. Three seed tiers are used throughout: $n=10$ paired seeds for
headline comparisons, $n=3$ as directional evidence for mechanism probes, and
single-seed pilots labelled as such; Table~\ref{tab:reporting-protocol} summarises
the metric hierarchy and reporting protocol applied uniformly across the chapter.
The results proceed from weakest to strongest stress: statistical heterogeneity,
the proximal coefficient, straggler asymmetry, learning-rate asymmetry and FedNova,
rare-class collapse, and mechanism checks.


% =============================================================================
% Table 5.1
% =============================================================================

\begin{table}[H]
\centering
\small
\caption{Reporting protocol and metric hierarchy. Applied uniformly across the
Results sections.}
\label{tab:reporting-protocol}
\begin{tabular}{p{0.20\linewidth} p{0.32\linewidth} p{0.40\linewidth}}
\hline
Component               & Choice                                       & Reason \\
\hline
Primary metric          & Test macro-F1 at best-val checkpoint         & Class-imbalanced; accuracy saturates at majority floor \\
Secondary metrics       & Per-class F1, per-class recall (confusion-matrix diagonal), rare-F1, rare-to-mel-nevi misrouting & Surface per-class failure that macro can mask; misrouting names the specific collapse mode \\
Rare-class definition   & $\{\text{dermato}, \text{melanoma}, \text{vascular}\}$ & Low support + high clinical priority (melanoma is mid-support but high-stakes) \\
$n = 10$ protocol       & Paired sign count + Wilcoxon where applicable & Headline statistical-heterogeneity comparisons \\
$n = 3$ protocol        & Directional only; no significance test       & Mechanism probes; power insufficient for $p$-values \\
$n = 1$ protocol        & Labelled ``pilot''; no headline use          & Heterogeneity ladder; single-class seed events possible \\
Runtimes                & Flower (central), pure-PyTorch (corroboration) & Flower carries full git+update-norm provenance \\
\hline
\end{tabular}
\end{table}




%%%%%%%%%%%%%%%%%%%%%%%%%%%%%%%%%%%%%%%%%%%%%%%%%%%%%%%%%%%%%%%%%%%%%%%%%%%%%%%
%% sec 5.2 -- Statistical heterogeneity alone gives limited optimiser separation
%%%%%%%%%%%%%%%%%%%%%%%%%%%%%%%%%%%%%%%%%%%%%%%%%%%%%%%%%%%%%%%%%%%%%%%%%%%%%%%

% =============================================================================
% Results sec 5.2 -- Statistical heterogeneity alone gives limited separation
% Compressed version. All numerics verified against VERIFICATION_SHEET.txt.
% Target ~500 words.
% =============================================================================

\section{Weak separation under statistical heterogeneity}
\label{sec:statistical-heterogeneity}

Under an IID partition (Flower, $n=10$ paired seeds), FedAvg and FedProx
differ by $\Delta = -0.007$ macro-F1 (FA $0.585$, FP $0.578$, FN $0.577$) ---
within seed noise. We treat this as a negative control: statistical
heterogeneity is a necessary condition for any FedAvg--FedProx gap on this
benchmark.

On the engineered balanced-paired seven-client partition we report three
estimates ($n=10$ paired seeds each): Flower paired $\Delta = +0.007$ (7/10
paired wins, full \texttt{git\_commit} and \texttt{client\_update\_norms}
provenance), Flower $\mu$-sweep peak at $\mu = 0.01$ $\Delta = +0.013$ (full
provenance, $\mu$ characterisation in \S\,\ref{sec:mu-sensitivity}), and
pure-PyTorch reference $\Delta = +0.027$ (9/10 wins, no \texttt{git\_commit},
no companion update-norm logs). All three measure the same partition through
slightly different code paths; they agree on sign and bracket the magnitude
at $0.01$--$0.03$. The Flower paired $\Delta = +0.007$ is the central
headline; the pure-PyTorch number is cross-runtime corroboration.

The pure-PyTorch $+0.027$ is concentrated in two minority classes --- melanoma
and actinic keratoses carry $61\%$ and $36\%$ of the macro-F1 advantage --- while
the other five classes, including the invariant mel-nevi, are flat
(Figure~\ref{fig:engineered-per-class}). The gain is the rebalancing of two
minority classes, not a uniform optimiser improvement.

Three further partitions replicate the small-magnitude finding: Dirichlet
$\alpha = 0.1$ and specialist 1-of-7 both give $\Delta = +0.008$ ($n=10$), and
the multi-seed node-pinned L4 control $\Delta = +0.005$ ($n=3$). The single-seed
heterogeneity-ladder pilot at L4 gives $\Delta = -0.030$, but that deficit is one
rare-class (vascular) outcome on one seed and the $n=3$ node-pinned promotion
averages it out. All Dirichlet aggregates use regenerated raw-JSON summaries; the
previously-committed analysis CSVs were stale and showed a spurious catastrophic
FedProx failure that the regenerated data do not support.

On the engineered partition, FedProx reduces test-macro-F1 SD from $0.025$
to $0.014$ ($-45\,\%$), reaches validation macro-F1 $= 0.40$ in $9.8$ vs
$12.7$ rounds ($1.3\times$ speed-up), and shows a paired mid-training peak
$\Delta = +0.094$ at round 50 that closes to $+0.007$ by round 150. These
trajectory characteristics are documented for the engineered partition only.

Statistical heterogeneity alone produces small, runtime-sensitive separation
(Table~\ref{tab:stat-het-headline}, Figure~\ref{fig:stat-het-forest}). The
next question is whether the proximal coefficient $\mu$ itself explains when
FedProx helps.


% =============================================================================
% Table 5.2 -- statistical heterogeneity and runtime sensitivity
% =============================================================================

\begin{table}[H]
\centering
\caption{Statistical heterogeneity and runtime sensitivity. Six cells:
IID, the engineered partition through two runtimes plus the $\mu$-peak,
Dirichlet $\alpha=0.1$, specialist 1-of-7, and the multi-seed node-pinned
L4 control. Dirichlet, specialist, and IID values from regenerated raw-JSON
aggregates.}
\label{tab:stat-het-headline}
\resizebox{\textwidth}{!}{%
\begin{tabular}{lccccccl}
\hline
Setting              & Runtime       & Clients & $n$ & FedAvg              & FedProx             & $\Delta$ & Note \\
\hline
IID                  & Flower        & 7       & 10  & $0.585 \pm 0.020$   & $0.578 \pm 0.020$   & $-0.007$ & negative control \\
Engineered (pure)    & pure-PyTorch  & 7       & 10  & $0.481 \pm 0.025$   & $0.508 \pm 0.014$   & $+0.027$ & 9/10 wins; partial provenance \\
Engineered (Flower)  & Flower        & 7       & 10  & $0.496$             & $0.503$             & $+0.007$ & 7/10 wins; \textbf{central} \\
Dirichlet $\alpha=0.1$ & Flower      & 7       & 10  & $0.480 \pm 0.036$   & $0.488 \pm 0.034$   & $+0.008$ & regenerated raw-JSON \\
Specialist 1-of-7    & Flower        & 7       & 10  & $0.437 \pm 0.024$   & $0.445 \pm 0.017$   & $+0.008$ & --- \\
Node-pinned L4       & Flower        & 2       & 3   & $0.492 \pm 0.019$   & $0.497 \pm 0.022$   & $+0.005$ & multi-seed L4 control \\
\hline
\end{tabular}%
}
\end{table}


% =============================================================================
% Figure 5.2(a) -- statistical heterogeneity forest plot
% =============================================================================

\begin{figure}[H]
\centering
\includegraphics[width=0.85\textwidth]{F_statistical_heterogeneity_forest.pdf}
\caption{Forest plot of $\Delta(\textrm{FedProx} - \textrm{FedAvg})$ macro-F1
across the six settings in Table~\ref{tab:stat-het-headline}. Whiskers show
$\pm 1$ sample SD across paired seeds; dashed line at $\Delta = 0$ marks no
separation. The IID negative control sits at the left; the five partitioned
settings cluster between $+0.005$ and $+0.027$.}
\label{fig:stat-het-forest}
\end{figure}


% =============================================================================
% Figure 5.2(b) -- engineered partition per-class trajectories
% =============================================================================

\begin{figure}[H]
\centering
\includegraphics[width=\textwidth]{F_engineered_per_class.pdf}
\caption{Per-class validation F1 trajectories on the engineered partition,
$n = 10$ paired seeds (FedAvg vs FedProx, mean $\pm 1$ SEM bands).
Two-by-four grid (seven classes plus legend). Melanoma and actinic-keratoses
panels show non-overlapping FedAvg/FedProx bands by mid-training --- the
two classes that carry the pure-PyTorch $+0.027$ macro-F1 advantage
($+0.114$ and $+0.067$ per-class F1, $61\,\%$ and $36\,\%$ of the share
respectively). The other five classes show overlapping bands throughout
150 rounds, including the dominant mel-nevi panel which is invariant near
its prevalence saturation.}
\label{fig:engineered-per-class}
\end{figure}



%%%%%%%%%%%%%%%%%%%%%%%%%%%%%%%%%%%%%%%%%%%%%%%%%%%%%%%%%%%%%%%%%%%%%%%%%%%%%%%
%% sec 5.3 -- mu sensitivity: the proximal term helps only in a narrow range
%%%%%%%%%%%%%%%%%%%%%%%%%%%%%%%%%%%%%%%%%%%%%%%%%%%%%%%%%%%%%%%%%%%%%%%%%%%%%%%

% =============================================================================
% Results sec 5.3 -- mu sensitivity
% Compressed version. Target ~300 words.
% =============================================================================

\section{Proximal coefficient: a narrow useful range}
\label{sec:mu-sensitivity}

If the small FedProx advantage of \S\,\ref{sec:statistical-heterogeneity} is
real, the natural next question is whether the proximal coefficient $\mu$ ---
the one knob that distinguishes FedProx from FedAvg --- governs it, and whether
a larger $\mu$ buys a larger advantage under heterogeneity. A five-cell sweep
of $\mu \in \{0.001, 0.01, 0.1, 1.0\}$ against the FedAvg
baseline ($\mu = 0$) on the engineered partition ($n=10$ paired seeds per
cell) answers both questions at once, tracing a clear inverted-$U$
(Table~\ref{tab:mu-sweep},
Figure~\ref{fig:mu-inverted-u}): peak near $\mu = 0.01$, dropping below
baseline by $\mu = 0.1$ and catastrophic at $\mu = 1.0$.

The peak is at $\mu = 0.01$ ($0.506$, $\Delta = +0.013$, 6/10 paired wins),
essentially tied with $\mu = 0.001$; by $\mu = 0.1$ FedProx is below the FedAvg
baseline ($0.493$), and at $\mu = 1.0$ it collapses to $0.419$
($\Delta = -0.074$, 0/10 wins --- every seed loses; Table~\ref{tab:mu-sweep}).
The $0.09$ macro-F1 spread between best and worst $\mu$ is an order of magnitude
larger than any partition-level $\Delta$ in
\S\,\ref{sec:statistical-heterogeneity}.

The $\mu$-sweep peak $\Delta = +0.013$ exceeds the
\S\,\ref{sec:statistical-heterogeneity} Flower paired $\Delta = +0.007$
because the sweep selects best-of-five on the same seeds; the two estimates
bracket the same small effect at $+0.005$ to $+0.015$.

The practical region of useful $\mu$ on this benchmark is a single decade
around $0.01$. The expectation that larger $\mu$ should be preferred under
stronger statistical heterogeneity is not supported here --- a benchmark
observation, not a refutation of the underlying analysis.

The $\mu = 1.0$ deficit concentrates in the rare and minority classes while the
dominant mel-nevi is essentially invariant (Table~\ref{tab:mu-per-class}) --- the
same rare-class shape as the dropped-straggler and asymmetric-LR failures
(\S\,\ref{sec:rare-class-collapse}), via a different mechanism
(over-regularisation rather than under-aggregation). Because $\mu$ alone is
sensitive and harmful when large, the much larger FedProx gains under straggler
protocols (\S\,\ref{sec:li-straggler}) cannot be attributed to the proximal
coefficient itself.


% =============================================================================
% Table 5.3 -- mu-sensitivity sweep
% =============================================================================

\begin{table}[H]
\centering
\caption{$\mu$ sensitivity on the engineered partition, $n=10$ paired
seeds per cell. ``Wins'' is paired comparisons where FedProx exceeds
FedAvg on the matched seed.}
\label{tab:mu-sweep}
\begin{tabular}{ccccl}
\hline
$\mu$    & macro-F1 (mean $\pm$ SD) & $\Delta$ vs FedAvg & Wins / 10 & Interpretation \\
\hline
$0$      & $0.493 \pm 0.020$        & ---               & ---       & FedAvg baseline                \\
$0.001$  & $0.504 \pm 0.021$        & $+0.011$          & 6         & near peak                      \\
$0.01$   & $\mathbf{0.506 \pm 0.017}$ & $\mathbf{+0.013}$ & 6       & \textbf{peak}                  \\
$0.1$    & $0.479 \pm 0.017$        & $-0.014$          & 2         & below baseline                 \\
$1.0$    & $0.419 \pm 0.013$        & $-0.074$          & 0         & catastrophic                   \\
\hline
\end{tabular}
\end{table}


% =============================================================================
% Figure 5.3 -- mu-sensitivity curve
% =============================================================================

\begin{figure}[H]
\centering
\includegraphics[width=0.85\textwidth]{F_mu_sensitivity_outcome_and_convergence.pdf}
\caption{Inverted-$U$ $\mu$ sensitivity curve on the engineered partition
($n=10$ per cell). $x$-axis: $\mu$ on $\log_{10}$ scale across
$\{0.001, 0.01, 0.1, 1.0\}$, FedAvg baseline as dashed reference line at
$0.493$. $y$-axis: test macro-F1 with $\pm 1$ SD whiskers.}
\label{fig:mu-inverted-u}
\end{figure}


% =============================================================================
% Table 5.3b -- per-class mu sensitivity (engineered partition)
% =============================================================================

\begin{table}[H]
\centering
\small
\caption{Per-class validation F1 mean across the $\mu$ sweep on the
engineered partition, $n = 10$ paired seeds per cell. SDs across seeds
range $0.003$ (mel-nevi) to $0.113$ (rare classes). ``$\Delta$ at $\mu=1.0$''
is the absolute per-class F1 deficit at the catastrophic cell relative to
the FedAvg baseline ($\mu = 0$). The deficit concentrates in rare and
minority classes; the dominant mel-nevi class is essentially invariant.
Source: \texttt{mu\_sensitivity\_flower/}.}
\label{tab:mu-per-class}
\resizebox{\textwidth}{!}{%
\begin{tabular}{lcccccc}
\hline
Class       & $\mu = 0$ & $\mu = 0.001$ & $\mu = 0.01$ & $\mu = 0.1$ & $\mu = 1.0$ & $\Delta$ at $\mu = 1.0$ \\
\hline
actinic     & $0.475$   & $0.470$       & $0.489$      & $0.434$     & $0.351$     & $-0.124$ \\
basal       & $0.497$   & $0.513$       & $0.497$      & $0.405$     & $0.438$     & $-0.060$ \\
b-kerat     & $0.405$   & $0.400$       & $0.423$      & $0.416$     & $0.328$     & $-0.077$ \\
dermato     & $0.259$   & $0.287$       & $0.283$      & $0.279$     & $0.158$     & $\mathbf{-0.101}$ \\
melanoma    & $0.244$   & $0.258$       & $0.304$      & $0.287$     & $0.216$     & $-0.028$ \\
mel-nevi    & $0.884$   & $0.885$       & $0.886$      & $0.886$     & $0.867$     & $-0.017$ \\
vascular    & $0.687$   & $0.713$       & $0.661$      & $0.646$     & $0.574$     & $\mathbf{-0.113}$ \\
\hline
\end{tabular}%
}
\end{table}



%%%%%%%%%%%%%%%%%%%%%%%%%%%%%%%%%%%%%%%%%%%%%%%%%%%%%%%%%%%%%%%%%%%%%%%%%%%%%%%
%% sec 5.4 -- Li-style straggler asymmetry (CENTREPIECE)
%%%%%%%%%%%%%%%%%%%%%%%%%%%%%%%%%%%%%%%%%%%%%%%%%%%%%%%%%%%%%%%%%%%%%%%%%%%%%%%

% =============================================================================
% Results sec 5.4 -- Centrepiece
% Compressed version. Target ~720 words (centrepiece stays substantive).
% =============================================================================

\section{Straggler asymmetry: \texorpdfstring{$\gamma$}{gamma}-inexact handling, not the proximal term}
\label{sec:li-straggler}

This section presents the central result of the thesis. Every preceding
section has shown the FedAvg--FedProx gap to be small; here it becomes
large, and a controlled four-condition decomposition attributes that gap
to its true cause --- not the proximal regulariser FedProx is named for,
but the $\gamma$-inexact partial-update handling it is conventionally
paired with. This is what turns the thesis from an optimiser comparison
into a regime map.

Before decomposing the mechanism, we first checked whether stragglers alone
create a FedProx advantage. They do not: in a 7-client IID random-straggler
control FedProx is slightly worse than FedAvg ($\Delta = -0.013$ macro-F1,
$n = 10$), but replacing the IID split with a heterogeneous one under the same
straggler protocol reverses the sign ($\Delta = +0.017$, $n = 10$). The advantage
appears only when slow or partially updated clients also carry non-redundant
class information; the next experiment isolates that mechanism in a two-client L4
decomposition with the rare-class client as the fixed straggler.

We next isolate the mechanism in a two-client L4 four-condition experiment,
with $n=3$ matched seeds (Table~\ref{tab:li-decomp-l4}). In this setting,
Client 1 contains the rare-class signal and is made the fixed straggler. The
four cells separate two effects that are otherwise bundled together in
FedProx-style straggler experiments: the proximal term and the decision to
accept or drop a straggler's partial update.

\begin{table}[H]
\centering
\caption{Li-style four-condition decomposition at L4 (two-client; $n=3$
matched seeds; mean $\pm$ sample SD across seeds).}
\label{tab:li-decomp-l4}
\small
\begin{tabular}{lcc}
\toprule
Condition                                       & Test macro-F1                & Rare-class F1 \\
\midrule
FedAvg, no straggler                            & $0.485 \pm 0.028$            & $0.316 \pm 0.060$ \\
FedAvg + drop                                   & $0.364 \pm 0.009$            & $0.000 \pm 0.000$ \\
FedProx + drop (control)                        & $0.368 \pm 0.005$            & $0.000 \pm 0.000$ \\
FedProx + $\gamma$-inexact                      & $\mathbf{0.479 \pm 0.010}$   & $\mathbf{0.291 \pm 0.018}$ \\
\bottomrule
\end{tabular}

\smallskip
\noindent\textit{Decomposition:} proximal-only $= 0.368 - 0.364 = +0.004$;
$\gamma$-inexact $= 0.479 - 0.368 = +0.111$; total $= +0.115$. Rare-F1 is
the mean over $\mathrm{RARE\_IDX} = \{\text{dermato}, \text{melanoma},
\text{vascular}\}$ and shows that the macro-F1 gap lives entirely in the
rare classes: both drop-handling cells flatline rare-F1 to exactly zero,
and only FedProx+$\gamma$-inexact rescues it.
\end{table}

The first two cells measure the cost of straggling: FedAvg falls from $0.485$
(no straggler) to $0.364$ when the partial update is dropped, so in L4 dropping
the straggler removes class information rather than merely slowing training. The
next two identify the cure. Adding the proximal term while still dropping changes
almost nothing (FedProx+drop $0.368$, $+0.004$ above FedAvg+drop), whereas
accepting the partial update recovers nearly the full baseline
(FedProx+$\gamma$-inexact $0.479$, a $+0.111$ gain over FedProx+drop). Of the
total $+0.115$ gain, then, about $96\%$ comes from $\gamma$-inexact acceptance and
only $+0.004$ from the proximal term. Because the design omits a
FedAvg+$\gamma$-inexact cell, this is evidence for the FedProx+$\gamma$-inexact
protocol, not proof that $\gamma$-inexact acceptance is universally beneficial on
its own.

Figure~\ref{fig:val-curves-extreme-gaps}~(A) shows the same result over
training: the FedAvg+drop and FedProx+drop trajectories overlap for all 150
rounds, while FedProx+$\gamma$-inexact separates from both around round 30 and
stabilises near $0.48$.

\begin{figure}[H]
\centering
\includegraphics[width=\textwidth]{F_val_curves_extreme_gaps.pdf}
\caption{Validation macro-F1 trajectories at the two settings with the
largest FedAvg vs FedProx differences (mean over three seeds; shaded
bands are $\pm 1$ sample SD). (A) The four-condition
decomposition at L4: the FedAvg+drop and FedProx+drop trajectories overlap
throughout the 150 rounds; only FedProx+$\gamma$-inexact separates and
recovers towards the no-straggler baseline ($\Delta$ up to $+0.115$).
(B) Perfect-storm L4 with $90\%$ random stragglers: FedAvg+drop flatlines
at its round-1 validation value, so the selected checkpoint is effectively
the untrained model, while FedProx+$\gamma$-inexact climbs to
$0.491$ at $\mu = 0.01$ ($\Delta$ up to $+0.404$ --- the largest
FedProx--FedAvg gap measured in the thesis).}
\label{fig:val-curves-extreme-gaps}
\end{figure}

Per-class trajectories (Figure~\ref{fig:l4-per-class-grid}) localise the gain to
the rare classes: mel-nevi stays near $0.89$ under all four conditions, while
dermatofibroma, melanoma, and vascular collapse under both drop cells and recover
only under FedProx+$\gamma$-inexact --- the mechanism is rare-class learning, not
a broad improvement.

\begin{figure}[H]
\centering
\includegraphics[width=\textwidth]{F_l4_four_condition_per_class_grid.pdf}
\caption{Per-class validation F1 trajectories under the four conditions of
the L4 decomposition (Table~\ref{tab:li-decomp-l4}); mean over three seeds.
The majority class (mel-nevi) stays near $0.89$ in every cell, so the
aggregate improvement is not driven by the dominant class. The three rare
classes (dermatofibroma, melanoma, vascular) collapse under both
drop-handling cells and recover most clearly under
FedProx+$\gamma$-inexact, identifying rare-class learning as the locus of
the mechanism.}
\label{fig:l4-per-class-grid}
\end{figure}

The perfect-storm experiment stresses the same ordering rather than re-deriving
it: on the same L4 setting with a $90\%$ per-round random-straggler rate (batch
$10$, no momentum), FedAvg+drop flatlines at its round-1 value --- best-val
selects the untrained init, mean test macro-F1 $0.087 \pm 0.049$ --- while
FedProx+$\gamma$-inexact recovers to $0.365 \pm 0.023$ at $\mu = 1.0$ and
$0.491 \pm 0.003$ at $\mu = 0.01$, the latter the largest FedProx--FedAvg gap in
the thesis ($\Delta = +0.404$; Figure~\ref{fig:val-curves-extreme-gaps}B). The
mechanism is the one the L4 decomposition isolated (Figure~\ref{fig:signal-flow-gamma}),
pushed to its limit; the smaller $\mu$ wins because looser anchoring leaves more
rare-class learning signal per round. This is a magnitude demonstration, not a
second causal decomposition: the schedule and arms differ from the four-condition
design, and FedAvg+drop collapses too completely for a clean
$\gamma$-inexact-vs-proximal split.

Finally, the L1 control tests whether partial-update acceptance helps simply
because it is a better generic straggler policy. It does not. L1 uses the
same two-client four-condition design, but the clients have identical label
proportions; the split changes quantity, not class information. In this
setting, all four cells are nearly flat (Table~\ref{tab:li-decomp-l1}):
FedAvg-no-straggler $0.601 \pm 0.010$, FedAvg+drop $0.599 \pm 0.006$,
FedProx+drop $0.593 \pm 0.022$, and FedProx+$\gamma$-inexact
$0.605 \pm 0.009$. Dropping the straggler costs almost nothing, and accepting
the partial update adds only $+0.012$, compared with $+0.111$ at L4.

\begin{table}[H]
\centering
\caption{L1 cross-partition control: same four-condition design as
Table~\ref{tab:li-decomp-l4} but on the L1 partition (quantity-only
heterogeneity, JS divergence $\approx 0$); $n=3$ matched seeds, mean
$\pm$ sample SD. The proximal-only and $\gamma$-inexact contributions both
collapse toward zero relative to L4, which is the cleanest negative
control for the rare-class-signal-rescue mechanism interpretation.}
\label{tab:li-decomp-l1}
\small
\begin{tabular}{lcc}
\toprule
Condition                                       & Test macro-F1        & Rare-class F1 \\
\midrule
FedAvg, no straggler                            & $0.601 \pm 0.010$    & $0.530 \pm 0.008$ \\
FedAvg + drop                                   & $0.599 \pm 0.006$    & $0.539 \pm 0.012$ \\
FedProx + drop (control)                        & $0.593 \pm 0.022$    & $0.537 \pm 0.039$ \\
FedProx + $\gamma$-inexact                      & $0.605 \pm 0.009$    & $0.556 \pm 0.014$ \\
\bottomrule
\end{tabular}

\smallskip
\noindent\textit{Decomposition (vs L4 in parentheses):}
proximal-only $= 0.593 - 0.599 = -0.006$ (L4: $+0.004$);
$\gamma$-inexact $= 0.605 - 0.593 = +0.012$ (L4: $+0.111$);
drop-cost $= 0.601 - 0.599 = +0.002$ (L4: $+0.121$). The $\gamma$-inexact
contribution shrinks by an order of magnitude when the dropped client
carries no class signal the other client lacks. The rare-F1 column makes
this concrete: at L1 every cell sits in $0.53$--$0.56$ (a $\pm 0.026$
band), while at L4 the two drop cells crash to exactly zero and only
$\gamma$-inexact recovers rare-F1 to $0.29$. There is nothing to rescue
when the dropped client carries no unique rare-class signal.
\end{table}

The L4/L1 contrast is therefore the key mechanism evidence. At L4, the
dropped client carries class information the other client does not have, so
accepting its partial update rescues rare-class learning. At L1, the dropped
client carries no unique class signal, so there is little to rescue. The
straggler mechanism is not generic drift control; it matters when system
heterogeneity suppresses statistically unique rare-class information.

The straggler experiments isolate one axis of system asymmetry.
Section~\ref{sec:lr-asymmetry} asks whether persistent learning-rate
imbalance --- a different axis --- changes the optimiser ranking.



%%%%%%%%%%%%%%%%%%%%%%%%%%%%%%%%%%%%%%%%%%%%%%%%%%%%%%%%%%%%%%%%%%%%%%%%%%%%%%%
%% sec 5.5 -- Learning-rate asymmetry reveals FedNova regime-dependence
%%%%%%%%%%%%%%%%%%%%%%%%%%%%%%%%%%%%%%%%%%%%%%%%%%%%%%%%%%%%%%%%%%%%%%%%%%%%%%%

% =============================================================================
% Results sec 5.5 -- Learning-rate asymmetry, FedNova regime-dependence
% =============================================================================

\section{Learning-rate asymmetry and FedNova regime-dependence}
\label{sec:lr-asymmetry}

This experiment tests a different form of system heterogeneity from
\S\,\ref{sec:li-straggler}. There, the rare-class client was slow and its
partial update could be dropped. Here, the rare-class client is not dropped:
it still participates in every round, but its learning rate is made smaller.
The question is whether an optimiser can preserve rare-class learning when
the rare client's update becomes progressively weaker.

The experiment is run on L4. Client 0 is the dominant client and keeps
$\textrm{lr}=0.01$ throughout. Client 1 holds the rare classes and is assigned
one of six smaller learning rates, producing asymmetry ratios from 1:1 to
50:1. Each ratio is tested with FedAvg, FedProx, and FedNova using $n=3$
matched seeds per cell, giving $54$ runs in total. Test macro-F1 across the
six-ratio envelope is reported in Table~\ref{tab:lr-envelope}.

\begin{table}[H]
\centering
\caption{Learning-rate-asymmetry envelope at L4, $n=3$ matched seeds per cell.
Client~0 (dominant) holds $\textrm{lr} = 0.01$ throughout; Client~1
(rare-class holder) takes the smaller rate that sets each ratio. Columns are
test macro-F1 for the three optimisers and the FedNova advantage over FedAvg
and FedProx. FedNova (bold) leads at every ratio and stays flat across the
envelope, while FedAvg and FedProx plateau at the rare-class floor by 5:1.}
\label{tab:lr-envelope}
\resizebox{\textwidth}{!}{%
\begin{tabular}{c|cc|ccc|cc}
\hline
Ratio & Client 0 lr & Client 1 lr & FedAvg & FedProx & FedNova & FN$-$FA & FN$-$FP \\
\hline
 1:1  & $0.0100$ & $0.0100$ & $0.457$ & $0.491$ & $\mathbf{0.509}$ & $+0.052$ & $+0.018$ \\
 2:1  & $0.0100$ & $0.0050$ & $0.425$ & $0.462$ & $\mathbf{0.526}$ & $+0.101$ & $+0.064$ \\
 5:1  & $0.0100$ & $0.0020$ & $0.375$ & $0.392$ & $\mathbf{0.504}$ & $+0.129$ & $+0.112$ \\
10:1  & $0.0100$ & $0.0010$ & $0.369$ & $0.369$ & $\mathbf{0.501}$ & $+0.132$ & $+0.132$ \\
20:1  & $0.0100$ & $0.0005$ & $0.367$ & $0.372$ & $\mathbf{0.486}$ & $+0.119$ & $+0.114$ \\
50:1  & $0.0100$ & $0.0002$ & $0.367$ & $0.367$ & $\mathbf{0.498}$ & $\mathbf{+0.131}$ & $\mathbf{+0.131}$ \\
\hline
\end{tabular}%
}
\end{table}

The pattern is clear (Table~\ref{tab:lr-envelope}). FedAvg and FedProx both
degrade as soon as Client~1's learning rate is reduced and plateau near $0.367$
macro-F1 by 5:1 --- the rare-class floor, where the model still learns the
dominant-client distribution but rare-class learning has largely collapsed.
FedNova behaves differently, holding $0.486$--$0.526$ across the full 1:1--50:1
envelope; even at 50:1 it reaches $0.498$ against $0.367$ for both FedAvg and
FedProx, a $+0.131$ macro-F1 gap. Rare-class F1 confirms this is not a macro-F1
artefact: it is near zero for FedAvg and FedProx by 10:1 (FA $0.016$, FP $0.010$)
but stays around $0.35$ for FedNova across the envelope
(Figure~\ref{fig:lr-asymmetry-envelope}).

\begin{figure}[H]
\centering
\includegraphics[width=\textwidth]{F_asymmetric_lr_L4.pdf}
\caption{Learning-rate asymmetry envelope at L4, $n=3$ per cell. Two
panels: test macro-F1 and rare-class F1 (mean over dermato, melanoma,
vascular) vs the Client 1 LR ratio on a $\log_{10}$ axis from 1:1 to 50:1.
FedAvg and FedProx plateau at the rare-class floor; FedNova stays flat
across the envelope.}
\label{fig:lr-asymmetry-envelope}
\end{figure}

The 5:1 cell is also where the chapter's metric choice earns its keep. Mean
test accuracy across the three matched seeds is $0.753$ for FedAvg and
$0.726$ for FedNova: ranking by raw accuracy would prefer FedAvg, even
though its macro-F1 ($0.375$) and rare-class F1 ($\approx 0.02$) make it
clear that FedAvg has lost the rare-class signal almost entirely while
FedNova has not (macro-F1 $0.504$, rare-F1 $\approx 0.37$). The accuracy
gap is small precisely because both models still get mel-nevi
($66.9\,\%$ of the test set) right; macro-F1 separates them because it
weights all seven classes equally. This is the operational reason \S\,\ref{sec:method-metrics}
ranks optimisers by macro-F1 and never by accuracy.

A simple interpretation is that FedAvg and FedProx allow the rare client's
update to become too small. When Client 1's learning rate is reduced, its
local model moves less during training, so the update it sends back to the
server carries weaker rare-class information. FedAvg does not correct this,
and the FedProx proximal term does not restore a rare-client update that has
already become too small. FedNova, by contrast, normalises client updates by
effective local work, which helps preserve Client 1's contribution in this
deterministic asymmetry setting.

The update-norm logs are mechanism-consistent: at 5:1 Client~1's update retains
$0.91$ of its 1:1 size under FedNova but only about half ($0.47$ FedAvg, $0.50$
FedProx), so the rare-client contribution is preserved under normalisation and
attenuated without it (\S\,\ref{sec:mechanism-robustness}).

FedNova's success here is not universal: it fails under random-$\tau$ straggler
heterogeneity, where the amount of local work varies randomly from round to
round. On the same 7-client engineered partition as the
\S\,\ref{sec:li-straggler} screening comparison (Regime~A, per-round random
schedule; \emph{not} an L4 experiment), FedNova obtains only
$0.237 \pm 0.089$ macro-F1, far below the matching FedAvg and FedProx cells
($n = 10$: FA $0.493 \pm 0.028$, FP $0.510 \pm 0.017$). Best-val selection masks
the severity: three of ten seeds already sit at the majority-class-only baseline
($\approx 0.114$, predicting mel-nevi for every image) at best-val, and all ten
reach it by the final round.

This contrast is the main point of the section. FedNova is robust when the
asymmetry is stable and deterministic: the rare client is consistently
weakened, and normalisation can compensate. It is brittle when the amount of
local work is random and noisy: normalisation can then amplify unstable
partial updates instead of correcting them. FedNova is therefore
regime-dependent, not universally superior
(Table~\ref{tab:fednova-regime}).

\begin{table}[H]
\centering
\small
\caption{FedNova is robust in one system-heterogeneity regime and brittle in
the other. Persistent, deterministic learning-rate asymmetry lets
$\tau$-normalisation preserve the weakened client's contribution; random
per-round $\tau$ --- the regime outside the stable-$\tau$ assumption of
\citet{wang2020fednova} --- causes the same mechanism to amplify noisy
partial updates and collapse to the majority-class baseline.}
\label{tab:fednova-regime}
\begin{tabular}{p{0.22\linewidth}p{0.32\linewidth}p{0.38\linewidth}}
\hline
FedNova regime & Result & Interpretation \\
\hline
Persistent LR asymmetry & macro-F1 holds $0.49$--$0.53$ across 1:1--50:1 & $\tau$-normalisation preserves smaller client's contribution \\
Random-$\tau$ stragglers & macro-F1 $= 0.237 \pm 0.089$; 10/10 final-round majority collapse (3/10 even at best-val) & $\tau$-normalisation amplifies noisy partial updates (outside Wang 2020) \\
\hline
\end{tabular}
\end{table}

Figure~\ref{fig:heterogeneity-escalation} places the large FedAvg--FedProx
gaps from the previous section on one axis. The progression from IID to
symmetric L4 to Li-style straggler asymmetry and then to perfect-storm shows
that optimiser gaps become large only in specific stress regimes. Together
with the FedNova result above, this reinforces the chapter's central point:
the optimiser ranking depends on the heterogeneity regime, not on the
optimiser name alone.

\begin{figure}[H]
\centering
\includegraphics[width=0.85\textwidth]{F_heterogeneity_escalation.pdf}
\caption{Heterogeneity-escalation synthesis. Within-pair
$\Delta(\textrm{FedProx} - \textrm{FedAvg})$ macro-F1 across five regime
cells of increasing severity: IID Flower negative control ($-0.007$),
symmetric L4 multi-seed node-pinned ($+0.005$), Li 2020 four-condition
contrast (FA+drop vs FP+$\gamma$-inex, $+0.115$), perfect-storm
$\mu = 1.0$ ($+0.278$), perfect-storm $\mu = 0.01$ ($+0.404$). The first
two cells sit inside seed noise; the last three are large,
mechanism-isolable, and regime-driven.}
\label{fig:heterogeneity-escalation}
\end{figure}

The weak protocols fail in a clinically specific way: they lose rare-class
sensitivity. The next section shows that this failure usually appears as
rare-class images being routed into the dominant mel-nevi class.



%%%%%%%%%%%%%%%%%%%%%%%%%%%%%%%%%%%%%%%%%%%%%%%%%%%%%%%%%%%%%%%%%%%%%%%%%%%%%%%
%% sec 5.6 -- Rare-class collapse into mel-nevi
%%%%%%%%%%%%%%%%%%%%%%%%%%%%%%%%%%%%%%%%%%%%%%%%%%%%%%%%%%%%%%%%%%%%%%%%%%%%%%%

% =============================================================================
% Results sec 5.6 -- Rare-class collapse + loss-side correction
% Compressed version. Target ~470 words.
% =============================================================================

\section{Rare-class collapse and the loss-side comparison}
\label{sec:rare-class-collapse}


\paragraph{Per-class collapse.}

The low-macro-F1 cells in \S\S\,\ref{sec:li-straggler}--\ref{sec:lr-asymmetry}
fail the same way: rare-class images are routed to mel-nevi
(Figure~\ref{fig:cross-experiment-per-class}). Mean rare-to-mel-nevi
misrouting (over dermato, melanoma, vascular) ranges from $24\,\%$ in the
best protocol (perfect-storm FP $\mu = 0.01$, $n=3$) to $67\,\%$ in the
worst (perfect-storm FA+drop, $n=3$), and $28\,\%$ even under the IID
Flower negative control ($n=10$) --- a class-imbalance saturation that the
FL protocol preserves or amplifies but does not itself create.

The D1 5:1 cell isolates the optimiser dimension --- melanoma per-class F1 is
$1.20\,\%$ (FedAvg), $6.13\,\%$ (FedProx), and $46.34\,\%$ (FedNova), the same
disease across three detection regimes spanning two orders of magnitude. Because
macro-F1 weights all seven classes equally, losing one class moves it by
$\approx 0.057$ --- the order of the entire engineered $\Delta$ in
\S\,\ref{sec:statistical-heterogeneity}. The same rare-class shape recurs across
aggregation dropout, learning-rate asymmetry, and the $\mu = 1.0$
over-regularisation cell (\S\,\ref{sec:mu-sensitivity}, Table~\ref{tab:mu-per-class}).

\begin{table}[H]
\centering
\small
\caption{Rare-class collapse across four representative cells, spanning the
straggler and learning-rate-asymmetry axes ($n=3$ each). Rare-F1, melanoma
per-class F1, and the rare-to-mel-nevi misrouting rate move together: where
rare signal reaches the global model (FedProx+$\gamma$-inexact, FedNova under
deterministic asymmetry) rare-F1 is preserved and misrouting is low; where it
does not (FedAvg+drop, FedAvg under 5:1 asymmetry) rare-F1 flatlines and
misrouting rises to $51$--$67\,\%$.}
\label{tab:collapse-cells}
\begin{tabular}{p{0.16\linewidth} p{0.14\linewidth} c c c p{0.32\linewidth}}
\hline
Protocol                  & Algorithm & Rare-F1 & Melanoma F1 & Misrouting & Interpretation \\
\hline
Perfect-storm + drop      & FedAvg    & $0.000$ & $0.00\,\%$  & $67\,\%$   & Rare client dropped; flatline \\
Perfect-storm + $\gamma$  & FedProx ($\mu=0.01$) & $0.357$ & $38.27\,\%$ & $24\,\%$ & $\gamma$-inex preserves rare signal \\
D1 5:1 LR asymmetry       & FedAvg    & $0.036$ & $1.20\,\%$  & $51\,\%$   & Rare client under-weighted; gradual erosion \\
D1 5:1 LR asymmetry       & FedNova   & $0.359$ & $46.34\,\%$ & $26\,\%$   & $\tau$-normalisation preserves contribution \\
\hline
\end{tabular}
\end{table}

\noindent
These four cells (Table~\ref{tab:collapse-cells}) span instant-collapse and
gradual-erosion modes across straggler and LR-asymmetry axes; the misrouting
rate co-varies with rare-F1 across all four.


\paragraph{Collapse modes.}

The collapse takes two temporal forms of one mechanism: an \emph{instant}
collapse when updates are dropped wholesale (perfect-storm FedAvg+drop flatlines
at its round-1 value, mean $0.087 \pm 0.049$), and a \emph{gradual} erosion when
they are merely under-weighted (D1 5:1 drifts $0.457 \to 0.375$ with rare-F1
falling in lockstep).


\paragraph{Loss-side imbalance correction (weighted CE and focal loss).}

To test whether the FedProx advantage might be class-imbalance correction
in disguise, we ran FedAvg and FedProx at L4 with three loss functions
($n=3$ per cell): plain cross-entropy, class-weighted CE with
inverse-frequency weights, and focal loss ($\gamma = 2$). The results are
in Table~\ref{tab:loss-side}.

\begin{table}[H]
\centering
\caption{Loss-side imbalance correction at L4 ($n=3$ per cell). Cells are
test macro-F1 $\pm$ sample SD with rare-class F1 in parentheses. The
strongest cell (bold) is FedAvg with weighted CE, which matches or exceeds
FedProx with plain CE: a direct loss-level correction reproduces the
FedProx-regime gain on this partition, so the FedProx advantage is not a
disguised class-imbalance correction. FedProx adds little on top of either
loss-side intervention.}
\label{tab:loss-side}
\begin{tabular}{lccc}
\hline
Algorithm & CE & Weighted CE & Focal \\
\hline
FedAvg  & $0.477 \pm 0.002$ ($0.285$) & $\mathbf{0.505 \pm 0.023}$ ($\mathbf{0.367}$) & $0.467 \pm 0.002$ ($0.317$) \\
FedProx & $0.478 \pm 0.016$ ($0.291$) & $0.503 \pm 0.014$ ($0.346$)                   & $0.476 \pm 0.022$ ($0.304$) \\
\hline
\end{tabular}
\end{table}

Three points follow (Table~\ref{tab:loss-side}). Weighted CE is the strongest
loss-side correction: FedAvg+wCE ($0.505$) marginally exceeds FedProx+CE
($0.478$), the CE-to-wCE step alone raising FedAvg rare-F1 by $+0.082$. Focal
trades macro for rare-recall (it raises rare-F1 but lowers macro at $\gamma=2$).
And FedProx adds little on top of either loss --- FP+wCE sits inside FA+wCE's
seed SD --- so FedProx acts as a drift/protocol intervention
(\S\,\ref{sec:li-straggler}), not a substitute for direct loss-side correction;
the two are combinable in principle but tested only separately here ($n=3$, L4
symmetric).

This does not conflict with the \S\,\ref{sec:li-straggler} result:
$\gamma$-inexact acceptance restores aggregation under straggler-induced signal
loss, while weighted CE re-weights rare classes during local training --- they
act at different pipeline stages (signal loss vs label saturation), and are
combinable in principle but tested only separately here.

The rare-class collapse documented here is clinically \emph{motivated},
not clinically \emph{validated}: DermaMNIST is a $28 \times 28$-pixel
benchmark and not a deployment artefact. The next section turns to the
per-round update-norm activity behind these per-class outcomes.


% =============================================================================
% Figure 5.7 -- cross-experiment per-class diagonal accuracy
% =============================================================================

\begin{figure}[H]
\centering
\includegraphics[width=\textwidth]{F_cross_experiment_per_class.pdf}
\caption{Per-class diagonal accuracy (correct-prediction rate per class)
across twelve protocol cells spanning baselines, Li 2020 straggler
conditions, perfect-storm, asymmetric LR (D1 1:1 and 5:1), and node-pinned
L4. The mel-nevi column saturates near $0.85$--$0.95$ in nearly every cell
while the dermato, melanoma, and vascular columns collapse to near zero in
the failing cells.}
\label{fig:cross-experiment-per-class}
\end{figure}



%%%%%%%%%%%%%%%%%%%%%%%%%%%%%%%%%%%%%%%%%%%%%%%%%%%%%%%%%%%%%%%%%%%%%%%%%%%%%%%
%% sec 5.7 -- Mechanism and robustness checks
%%%%%%%%%%%%%%%%%%%%%%%%%%%%%%%%%%%%%%%%%%%%%%%%%%%%%%%%%%%%%%%%%%%%%%%%%%%%%%%

% =============================================================================
% Results sec 5.7 -- Mechanism and robustness checks
% Compressed version. Target ~390 words.
% =============================================================================

\section{Mechanism and robustness}
\label{sec:mechanism-robustness}

\paragraph{Update-norm activity.}

Update norms ($\lVert \Delta w \rVert$, the size of the client update returned to
the server) are used here only as a supporting diagnostic, not as causal proof.
Two facts emerge. First, FedProx does what its proximal term intends: on the
engineered partition ($n=10$) the mean per-round update norm falls from $1.523$
(FedAvg) to $1.029$ (FedProx, $\mu=0.01$), a $32\%$ reduction
(Figure~\ref{fig:update-norms}). This is a local-training effect and does not
contradict \S\,\ref{sec:li-straggler}, where the straggler-regime gain came from
accepting the rare client's partial update rather than from the proximal term.

Second, update size tracks rare-class performance only within a controlled
protocol. Pooled across 78 L4 runs, larger dominant-client update norms
correlate with higher rare-class F1 (Spearman $r=+0.634$), but the
Client~1-to-Client~0 update-size ratio does not ($r=-0.045$) because the pool
mixes mechanisms (drop, $\gamma$-inexact, LR, $\mu$, loss). Inside the
asymmetric-LR experiment alone, where only Client~1's learning rate varies,
rare-class F1 closely follows the rare client's relative update strength
($r=+0.80$ FedAvg, $r=+0.87$ FedProx; cf.\ \S\,\ref{sec:lr-asymmetry}). These are
mechanism-consistent correlations, not causal proof; the main claims rest on the
controlled contrasts in \S\,\ref{sec:li-straggler} and \S\,\ref{sec:lr-asymmetry}.

\paragraph{Protocol-flip audit.}

The protocol-flip audit checks whether the algorithm ranking depends on the
reporting checkpoint (best-validation vs final round vs last-$K$ mean); a flip is
a sign change of the paired difference across these choices. Large effects do not
flip and small ones do: $0/9$ flips in the large-effect settings (Li-style
decomposition, perfect-storm) versus $3/9$ in the small-effect settings. The
centrepiece \S\,\ref{sec:li-straggler} effect ($+0.115$) and the
\S\,\ref{sec:lr-asymmetry} FedNova gap ($+0.131$) therefore survive
checkpoint choice, while the small \S\,\ref{sec:statistical-heterogeneity}
separations should be read as fragile. Further diagnostics (the
$B$-dissimilarity proxy, training-instability metric, asymmetric-$\mu$ Yao test,
early-warning collapse detector, and small-hospital case study) were also
examined; none is load-bearing for the main conclusions and they are omitted for
brevity. Both checks support the same message: large regime-dependent effects are
robust, small symmetric-protocol differences are not.

% =============================================================================
% Figure 5.8 -- update-norm mechanism
% =============================================================================

\begin{figure}[H]
\centering
\includegraphics[width=0.85\textwidth]{F_update_norms.pdf}
\caption{Update-norm mechanism on the engineered partition. Mean per-round
$\lVert \Delta w \rVert$ trajectories per client for FedAvg vs FedProx
($n=10$ seeds); FedProx tracks lower throughout, consistent with the
proximal term's drift-control intent. Cross-protocol scatter panel
(78 L4 runs across 8 algorithm $\times$ protocol cells): Spearman
$r = +0.634$ between dominant-client mean update-norm and rare-class F1;
relative-influence ratio $r = -0.045$ across protocols.}
\label{fig:update-norms}
\end{figure}



%%%%%%%%%%%%%%%%%%%%%%%%%%%%%%%%%%%%%%%%%%%%%%%%%%%%%%%%%%%%%%%%%%%%%%%%%%%%%%%
%% sec 5.8 -- Integrated summary of findings
%%%%%%%%%%%%%%%%%%%%%%%%%%%%%%%%%%%%%%%%%%%%%%%%%%%%%%%%%%%%%%%%%%%%%%%%%%%%%%%

% =============================================================================
% Results sec 5.8 -- Integrated summary of findings
% Compressed. Target ~350 words. Closes with the central thesis sentence
% verbatim per guide sec 1.
% =============================================================================

%% [Batch 1] Integrated-summary synthesis and the claim--evidence--caveat table
%% (tab:final-claims) were relocated into the Discussion (\S4.1). The duplicated
%% signal-flow restatement was dropped per the restructuring plan.



%%%%%%%%%%%%%%%%%%%%%%%%%%%%%%%%%%%%%%%%%%%%%%%%%%%%%%%%%%%%%%%%%%%%%%%%%%%%%%%
%% sec 5.9 -- Limitations
%%%%%%%%%%%%%%%%%%%%%%%%%%%%%%%%%%%%%%%%%%%%%%%%%%%%%%%%%%%%%%%%%%%%%%%%%%%%%%%

% =============================================================================
% Results sec 5.9 -- Limitations
% Compressed. Target ~200 words. One-line bullets, terse.
% =============================================================================

%% [Batch 1] The Limitations list was relocated into the Discussion and
%% Conclusion chapter (\S4.2 Limitations).


%% =============================================================================
%% END OF RESULTS BODY
%% =============================================================================


%% =============================================================================
%% CHAPTER 6: DISCUSSION AND CONCLUSION
%% =============================================================================
\chapter{Discussion and Conclusion}
\label{ch:discussion}

\section{Discussion}
\label{sec:disc}

The results are best read not as a contest between three optimisers but as a map
of when each preserves rare-class signal in the global model. Optimiser
performance is regime-dependent for one reason: rare-class performance depends on
whether the rare client's contribution survives aggregation, and different
optimisers protect or forfeit it under different heterogeneity. One rule organises
the six experiments --- rare-class F1 is preserved precisely when rare-client
signal continues to reach the global model. Figure~\ref{fig:regime-map-summary}
summarises the resulting map across the regimes studied.

\begin{figure}[t]
\centering
\includegraphics[width=\textwidth]{F_regime_map_summary.pdf}
\caption{Regime map: a qualitative synthesis of Chapter~\ref{ch:results}, not new
results. Each cell records the observed outcome for an optimiser under one
heterogeneity regime --- preserved / competitive / rescue / robust (green),
fragile or no extra gain (amber), collapse or brittle (red), or not tested in this
thesis (grey, used wherever an optimiser was not run in that regime). The map
records only the regimes studied here, on a single $28 \times 28$ medical
benchmark with simulated statistical and system heterogeneity; it is not a general
law of federated learning and implies no universal optimiser ranking. Mechanism
probes are $n = 3$ (directional); Table~\ref{tab:final-claims} gives the per-claim
evidence, strength labels, and caveats.}
\label{fig:regime-map-summary}
\end{figure}

Read through that rule, the findings cohere. Statistical heterogeneity alone
separates FedAvg and FedProx only slightly and runtime-sensitively
(\S\ref{sec:statistical-heterogeneity}), too fragile to justify a universal
preference, and the proximal coefficient is sharply tuned: useful near
$\mu = 0.01$ but harmful at $\mu = 1.0$, where it damages exactly the rare classes
(\S\ref{sec:mu-sensitivity}). This already indicates that the proximal term acts
on local drift, not on whether a client's signal is admitted. The L4/L1
decomposition makes the point precise (\S\ref{sec:li-straggler}): under straggler
asymmetry the FedProx advantage of about $+0.115$ macro-F1 comes almost entirely
($+0.111$) from $\gamma$-inexact acceptance of the straggler's partial update and
only $+0.004$ from the proximal term, while the matched L1 control --- where the
suppressed client carries nothing unique --- collapses the effect, confirming that
admitting the partial update matters only when that client is the sole source of
rare-class signal. FedNova is the same rule from another angle
(\S\ref{sec:lr-asymmetry}): work-normalisation restores a consistently weakened
client under deterministic learning-rate asymmetry (an advantage of about
$+0.131$ macro-F1 at the most extreme ratio) but amplifies noise and collapses to
roughly $0.237$ under random, round-to-round straggling. Finally, class-weighted
cross-entropy matches or exceeds FedProx with plain cross-entropy on the symmetric
partition (\S\ref{sec:rare-class-collapse}), locating imbalance correction as a
loss-side layer distinct from the aggregation-and-protocol layer where FedProx
acts; the two address different failure modes and could in principle be combined.

These findings refine rather than contradict prior work. FedAvg
\citep{mcmahan2017fedavg} behaves as expected when clients are well matched; for
FedProx \citep{li2020fedprox}, the decomposition indicates --- within this
benchmark and this simulated protocol --- that the straggler-regime benefit is
carried mainly by inexact-update acceptance rather than by the proximal term, an
attribution to the FedProx-plus-$\gamma$-inexact protocol as configured here
rather than a claim that the method underperforms; and for FedNova
\citep{wang2020fednova} the random-$\tau$ collapse is consistent with operating
outside its stable-work assumption, not with a flaw in its analysis. No result
here should be read as refuting these analyses. The clinical reading stays
cautious too: rare lesions collapsing into melanocytic nevi is the shared failure
mode, but DermaMNIST is a $28 \times 28$ benchmark, so this is clinically
motivated, not validated.

The five headline findings, with explicit strength labels and caveats, are
collected in Table~\ref{tab:final-claims}: the small, runtime-sensitive
statistical-heterogeneity separation (Flower $\Delta = +0.007$, IID control
$-0.007$); the inverted-$U$ $\mu$ sensitivity (catastrophic $\Delta = -0.074$ at
$\mu = 1.0$); the L4 decomposition in which $\gamma$-inexact handling carries
$+0.111$ of the $+0.115$ straggler-regime gain while the proximal term carries
only $+0.004$, with the L1 control flattening it to $+0.012$; FedNova's
regime-dependence ($+0.131$ at 50:1 deterministic asymmetry versus collapse to
$0.237$ under random-$\tau$); and rare-class collapse into mel-nevi ($24$--$67\%$
misrouting), against which weighted CE on FedAvg ($0.505$) matches FedProx+CE
($0.478$).

Federated optimiser performance in imbalanced medical FL is regime-dependent:
FedAvg remains competitive under IID or symmetric low-risk protocols;
FedProx becomes valuable under Li-style straggler asymmetry mainly because
$\gamma$-inexact partial-update handling preserves the rare-client signal;
FedNova is robust under the tested learning-rate asymmetry but fails under
random-$\tau$ straggler regimes; and across failing settings, the common
clinical failure mode is rare-class collapse into the majority class
(mel-nevi). These results do not establish a universal ranking of FedAvg,
FedProx, and FedNova. They instead map which optimiser/protocol choices
preserve or lose rare-class performance under different heterogeneity
regimes.

\begin{table}[H]
\centering
\small
\caption{Final claim--evidence--caveat summary. Strength labels: \emph{strong}
= multi-seed, mechanism-isolable, checkpoint-robust; \emph{medium} =
multi-seed but small effect or single-protocol; \emph{exploratory} =
correlational, single-protocol, or $n=3$ without mechanism isolation.}
\label{tab:final-claims}
\begin{tabular}{p{3.0cm}p{3.6cm}p{1.6cm}p{2.6cm}p{2.6cm}}
\hline
Claim & Evidence & Strength & Caveat & Implication \\
\hline
FedProx not free under IID & Flower IID $\Delta = -0.007$ ($n=10$) & strong & IID is artificial & Reach for FedProx only when heterogeneity exists \\
Stat-het alone gives fragile separation & Flower engineered $\Delta = +0.007$ (7/10 wins, $n=10$) & medium & Runtime-sensitive (PyTorch $+0.027$) & Do not recommend FedProx on partition alone \\
$\mu$ must be tuned & Sweep $n=10$: peak $+0.013$ at $\mu = 0.01$; $-0.074$ at $\mu = 1.0$ & strong & Engineered partition only & Default $\mu = 0.01$; do not raise under heterogeneity \\
$\gamma$-inexact dominates straggler gain & L4 four-condition $n=3$: total $+0.115$, $\gamma$-inex $+0.111$; L1 control $+0.012$ & medium--strong & Simulated stragglers; no FedAvg+$\gamma$-inex standalone cell & Use FedProx + $\gamma$-inexact under straggler protocols \\
FedNova regime-dependent & LR envelope: FN holds $0.49$--$0.53$ across 1:1--50:1 ($n=3$); random-$\tau$ collapses to $0.237$ ($n=10$, 10/10 final-round majority collapse) & medium--strong & Wang 2020 stable-$\tau$ assumption violated by random-$\tau$ & Prefer FedNova under persistent LR asymmetry, not under random-$\tau$ \\
Rare-class collapse is the clinical failure mode & Misrouting $24\,\%$--$67\,\%$; D1 5:1 melanoma F1 $= 1.2$/$6.1$/$46.3\,\%$ ($n=3$) & strong & Single dataset, $28 \times 28$ resolution & Report per-class metrics; macro-F1 alone hides it \\
wCE competes with FedProx & L4 $n=3$: FA+wCE $0.505$ vs FP+CE $0.478$ & medium & $n=3$, L4 symmetric only & Loss-side rebalancing is a competitive intervention \\
Update-norms support drift-control reading & Engineered: FA $1.523$ vs FP $1.029$ ($-32\,\%$); cross-protocol $r = +0.634$ & exploratory & Within-protocol $r$ does not extrapolate & Drift-control reading is supported, not proven \\
\hline
\end{tabular}
\end{table}

\section{Limitations}
\label{sec:limitations}

Several limitations bound these conclusions. The mechanism probes use three
matched seeds and are reported as directional; stragglers are simulated through
local-epoch counts rather than wall-clock deadlines; the two runtimes disagree in
magnitude on the engineered partition, so only its direction is robust; the study
uses a single low-resolution dataset and mixes two federation sizes; it omits
variance-reduced and adaptive comparators such as SCAFFOLD
\citep{karimireddy2020scaffold} and adaptive federated optimisers
\citep{reddi2021adaptive}; and evaluation is by argmax predictions only. The
thesis therefore establishes neither a universal ranking nor a clinically
validated claim. The map still gives regime-indexed guidance: under near-IID
conditions FedAvg is competitive; under statistical heterogeneity alone the
FedProx gain is small and the proximal coefficient should be kept small; when
stragglers coincide with a client holding unique rare-class signal, FedProx with
$\gamma$-inexact acceptance is valuable; under deterministic learning-rate
asymmetry FedNova is promising but should not be assumed safe under random
straggling; and under severe imbalance, per-class metrics should be reported and
weighted cross-entropy considered as a competitive complement.

\begin{itemize}
\setlength\itemsep{0.15em}
\item Single dataset at low resolution (DermaMNIST $28 \times 28$); no
replication on Camelyon, full-resolution ISIC, or non-dermatology MedMNIST
subsets.
\item Mixed federation sizes: 7-client for the engineered headline and
statistical-heterogeneity settings; 2-client for the mechanism probes
(Li 2020 four-condition, perfect-storm, asymmetric LR, weighted-CE,
asymmetric $\mu$, node-pinned L4). No cross-scale generalisation claim is
made between the two sizes.
\item Local-epoch budget fixed at $E = 20$; straggler protocols are
simulated via local-epoch counts rather than wall-clock deadlines.
\item Seed counts vary across experiments ($n = 10$ for headline comparisons,
$n = 3$ for the mechanism and asymmetry probes --- Li 2020 four-condition,
perfect-storm, asymmetric LR, weighted-CE, asymmetric $\mu$ --- and single-seed
pilots labelled as such) because HPC access was restricted. This is reported
transparently per experiment and should be weighed when interpreting the
uncertainty of the $n = 3$ claims, for which no significance test is applied.
\item The heterogeneity ladder is a single-seed pilot and is reported as
such; multi-seed promotion exists only at L4 ($n = 3$).
\item Runtime divergence between pure-PyTorch ($+0.027$) and Flower
($+0.007$) on the engineered partition limits magnitude claims; only
direction is robust across runtimes.
\item No SCAFFOLD, FedAdam, or FedYogi comparator; the optimiser space is
FedAvg, FedProx, FedNova only.
\item No AUC or calibration metrics; predictions are argmax only.
\item Per-client validation distributions were not logged centrally;
q-FFL-style fairness CDFs and Sheller-style cross-client generalisation
grids cannot be produced from existing artefacts.
\end{itemize}

\section{Conclusion}
\label{sec:conc}

This thesis asked not whether a federated optimiser is best for imbalanced medical
image classification but \emph{when} each preserves rare-class performance. The
answer is a regime map governed by a single rule: rare-class performance is
preserved when rare-client signal reaches the global model. Statistical
heterogeneity alone yields only a fragile FedProx--FedAvg separation (RQ1); under
straggler asymmetry the FedProx advantage is carried mainly by $\gamma$-inexact
partial-update handling rather than the proximal term, and only when the
suppressed client holds unique rare-class signal (RQ2); FedNova is
regime-dependent, robust under deterministic learning-rate asymmetry but brittle
under random stragglers (RQ3); and rare-class collapse into melanocytic nevi is
the common failure mode, with weighted cross-entropy showing the FedProx advantage
is not class-imbalance correction in disguise (RQ4). Concretely, in the
perfect-storm stress this margin is large: FedProx with $\gamma$-inexact
acceptance reaches $0.491$ test macro-F1 where FedAvg with straggler dropping
collapses to $0.087$ --- the thesis's largest FedAvg--FedProx gap, $\Delta =
+0.404$. Future work should combine
$\gamma$-inexact aggregation with a weighted-cross-entropy local objective, which
were tested only separately here; validate the map on higher-resolution and
additional datasets with real wall-clock straggler protocols; add variance-reduced
and adaptive optimisers; and log per-client validation to quantify fairness across
sites. In imbalanced medical federated learning, the pertinent question is not
which optimiser to use but which optimiser and protocol preserve rare-class signal
under the heterogeneity at hand: optimiser choice is regime-dependent, and it is
the regime, not the method's name, that should drive the decision.


%% =============================================================================
%% BIBLIOGRAPHY
%% =============================================================================
\bibliographystyle{plainnat}
\bibliography{references}


%% =============================================================================
%% APPENDICES
%% =============================================================================
\appendix

\chapter{Provenance, reproducibility, and supplementary diagnostics}
\label{app:repro}

\section{Runtimes and implementation provenance}
\label{app:provenance}

The implementation modules referenced abstractly in the Methods chapter are:
aggregation and the pure-PyTorch loop under \texttt{fl/} (\texttt{aggregation.py},
\texttt{server\_loop.py}, \texttt{local\_train.py}, \texttt{system\_het.py},
\texttt{class\_imbalance.py}, \texttt{evaluation.py}), the partitioner under
\texttt{data/partition.py}, the Flower strategies under \texttt{fl\_flower/}
(\texttt{strategy\_fednova.py}, \texttt{strategy\_straggler\_dropping.py}), and the
guards under \texttt{tests/} covering the proximal term, the
$\mu = 0 \Leftrightarrow$ FedAvg equivalence, the FedNova normaliser, straggler
dropping, partition determinism, and framework-provenance fields.

Flower production runs save \texttt{history\_*.csv} (per-round validation),
\texttt{test\_at\_best\_*.json} (selected-round metrics plus a provenance
dictionary --- git commit, hostname, framework version, Python/PyTorch versions,
ISO start/finish timestamps), \texttt{test\_predictions\_*.npz} (argmax
predictions and labels, used for confusion matrices and rare-to-mel-nevi
misrouting in §\ref{sec:rare-class-collapse}), and optional
\texttt{client\_update\_norms\_*.csv} (diagnostic only). Older pure-PyTorch
headline runs retain \texttt{history\_*.csv} and \texttt{test\_at\_best\_*.json}
but lack full git-commit provenance, prediction NPZs, and companion update-norm
logs, which is why the chapter treats them as direction-only corroboration. Every
numerical claim in Chapter~\ref{ch:results} traces to raw artefacts or freshly
regenerated summaries, never to stale committed CSVs.

\end{document}
